\chapter{表示论入门}

表示论的基本思想是：与其直接研究一个抽象群 $G$，不如研究它在线性空间上的作用。这样一来，群论中的乘法规律就被翻译成矩阵乘法规律；而矩阵、线性映射、特征值、迹、内积等线性代数工具便都可以投入使用。对于模形式理论而言，表示论提供了两条极其重要的桥梁：其一，Hecke 算子族构成交换代数，它们在模形式空间上的作用可以用表示论语言统一描述；其二，$\GL_2$ 及其子群的表示正是现代自守表示论的起点。本章的目标是给出有限群表示论的入门框架，并在结尾说明它如何自然地通向 Hecke 理论。

本章默认群 $G$ 为有限群；除特别说明外，表示空间取有限维向量空间，系数域记为 $F$。在讨论特征标时，我们主要取 $F=\C$。

\section{群表示的定义}

\begin{definition}
设 $G$ 是一个群，$V$ 是域 $F$ 上的有限维向量空间。一个\emph{$G$ 在 $V$ 上的线性表示}，或简称\emph{$G$ 的一个表示}，是一个群同态
\[
\rho:G\longrightarrow \GL(V).
\]
这里 $\GL(V)$ 表示 $V$ 的所有可逆线性变换组成的群。向量空间 $V$ 称为该表示的\emph{表示空间}，其维数 $\dim_F V$ 称为表示的\emph{次数}（degree）。
\end{definition}

换言之，每个 $g\in G$ 都对应一个可逆线性变换 $\rho(g):V\to V$，并满足
\[
\rho(1)=\Id_V,\qquad \rho(gh)=\rho(g)\rho(h)\quad (g,h\in G).
\]
为了简化记号，我们常把 $\rho(g)(v)$ 写成 $g\cdot v$。于是表示就是一个满足
\[
1\cdot v=v,\qquad g\cdot(h\cdot v)=(gh)\cdot v
\]
的线性群作用。

\begin{example}[平凡表示]
任意群 $G$ 在一维空间 $F$ 上都有一个最简单的表示：令
\[
\rho(g)=\Id_F\qquad (g\in G).
\]
这称为 $G$ 的\emph{平凡表示}，记作 $\mathbf{1}_G$ 或简记为 $\mathbf{1}$。它的次数为 $1$。
\end{example}

\begin{example}[正则表示]
设 $F[G]$ 是以 $\{e_g:g\in G\}$ 为基的 $F$-向量空间，其中基向量由群元素标记。定义
\[
\rho_{\mathrm{reg}}(x)(e_g)=e_{xg}\qquad (x,g\in G).
\]
则 $\rho_{\mathrm{reg}}:G\to \GL(F[G])$ 是一个表示，称为\emph{左正则表示}。它的次数是 $|G|$。

验证并不困难：
\[
\rho_{\mathrm{reg}}(xy)(e_g)=e_{xyg}=
ho_{\mathrm{reg}}(x)(e_{yg})
=
ho_{\mathrm{reg}}(x)\rho_{\mathrm{reg}}(y)(e_g),
\]
故 $\rho_{\mathrm{reg}}(xy)=\rho_{\mathrm{reg}}(x)\rho_{\mathrm{reg}}(y)$。
\end{example}

\begin{example}[置换表示]
设 $G$ 作用在有限集合 $X$ 上。令 $F[X]$ 为以 $\{e_x:x\in X\}$ 为基的向量空间。定义
\[
\rho_X(g)(e_x)=e_{g\cdot x}.
\]
则得到一个表示 $\rho_X:G\to \GL(F[X])$，称为与该群作用对应的\emph{置换表示}。它的次数是 $|X|$。

特别地，当 $G=S_n$ 作用在 $\{1,2,\dots,n\}$ 上时，便得到 $n$ 维置换表示。
\end{example}

\begin{example}[$S_n$ 的符号表示]
定义
\[
\sgn:S_n\longrightarrow \{\pm 1\}\subset F^{\times}
\]
为置换的奇偶性映射，即偶置换送到 $1$，奇置换送到 $-1$。因为
\[
\sgn(\sigma\tau)=\sgn(\sigma)\sgn(\tau),
\]
故它给出 $S_n$ 的一个一维表示，称为\emph{符号表示}。当 $\Char(F)\neq 2$ 时，它与平凡表示不同；当 $\Char(F)=2$ 时，$-1=1$，故它退化为平凡表示。
\end{example}

\begin{example}[循环群的二维表示]
设 $G=C_4=\langle r\mid r^4=1\rangle$，取 $F=\R$，定义
\[
\rho(r)=\mat{0&-1\\1&0}.
\]
这是平面上逆时针旋转 $90^\circ$ 的矩阵。因为 $\rho(r)^4=I_2$，故可唯一延拓为 $C_4$ 的一个二维表示。这说明几何对称经常天然地给出群表示。
\end{example}

给定一个表示后，选取 $V$ 的一组基，就可以把每个 $\rho(g)$ 写成矩阵，于是表示也可视为一个矩阵值群同态
\[
\rho:G\to \GL_n(F),\qquad n=\dim_F V.
\]
不同的基会给出相似的矩阵，因此“同一个表示”在不同基下的矩阵形式不应被视为本质不同。

\begin{definition}
设 $\rho:G\to \GL(V)$ 与 $\rho':G\to \GL(W)$ 是两个表示。如果存在线性同构
\[
T:V\longrightarrow W
\]
使得对一切 $g\in G$ 都有
\[
T\rho(g)=\rho'(g)T,
\]
则称 $\rho$ 与 $\rho'$ \emph{等价}，记作 $\rho\cong \rho'$。此时 $T$ 称为一个\emph{交织同构}（intertwining isomorphism）。
\end{definition}

上式也可写成
\[
\rho'(g)=T\rho(g)T^{-1},
\]
因此等价表示在矩阵语言中就是同时相似的矩阵组。

\begin{example}
仍以置换表示为例。设 $G$ 作用在两个有限集合 $X,Y$ 上，若存在双射 $\varphi:X\to Y$ 满足
\[
\varphi(g\cdot x)=g\cdot \varphi(x)
\qquad (g\in G,\ x\in X),
\]
则相应的置换表示 $F[X]$ 与 $F[Y]$ 等价。交织同构由 $e_x\mapsto e_{\varphi(x)}$ 给出。
\end{example}

\begin{proposition}
设 $\rho:G\to \GL(V)$ 是一个表示，$v_1,\dots,v_n$ 与 $w_1,\dots,w_n$ 是 $V$ 的两组基。设 $A_g$ 是 $\rho(g)$ 在基 $v_1,\dots,v_n$ 下的矩阵，$B_g$ 是它在基 $w_1,\dots,w_n$ 下的矩阵，则存在同一个可逆矩阵 $P\in \GL_n(F)$，使得
\[
B_g=PA_gP^{-1}\qquad (g\in G).
\]
反之，若两组矩阵值映射 $A,B:G\to \GL_n(F)$ 满足 $B_g=PA_gP^{-1}$，则它们给出等价表示。
\end{proposition}

\begin{proof}
令 $P$ 为从基 $v_1,\dots,v_n$ 到基 $w_1,\dots,w_n$ 的换基矩阵，则线性代数告诉我们同一线性变换在两组基下的矩阵满足相似关系
\[
B_g=PA_gP^{-1}.
\]
这对每个 $g\in G$ 同时成立。

反过来，若 $B_g=PA_gP^{-1}$，令 $T:F^n\to F^n$ 为矩阵 $P$ 所代表的线性同构，则
\[
TA_g=B_gT
\]
对每个 $g$ 成立，因此 $T$ 是交织同构，两表示等价。
\end{proof}

\begin{remark}
从表示论角度看，等价表示应被视为“同一个表示的不同坐标表达”。因此研究表示时，真正重要的是：哪些子空间在群作用下保持不变？表示能否分解成更简单的块？这些问题将在下一节开始讨论。
\end{remark}

\section{不可约表示与 Maschke 定理}

\begin{definition}
设 $\rho:G\to \GL(V)$ 是一个表示。若子空间 $W\subseteq V$ 满足对每个 $g\in G$ 都有
\[
\rho(g)(W)\subseteq W,
\]
则称 $W$ 是 $V$ 的一个\emph{子表示}或\emph{不变子空间}。

若 $V\neq 0$ 且它只有平凡的不变子空间 $0$ 和 $V$，则称 $V$（或表示 $\rho$）是\emph{不可约的}。否则称为\emph{可约的}。
\end{definition}

\begin{example}
平凡表示一定不可约，因为其表示空间是一维的。
\end{example}

\begin{example}
设 $S_3$ 在 $\C^3$ 上通过置换坐标作用：
\[
\sigma\cdot (x_1,x_2,x_3)=(x_{\sigma^{-1}(1)},x_{\sigma^{-1}(2)},x_{\sigma^{-1}(3)}).
\]
则
\[
U=\{(x,x,x):x\in \C\}
\]
是一个一维不变子空间，给出平凡子表示；
\[
W=\{(x_1,x_2,x_3)\in \C^3:x_1+x_2+x_3=0\}
\]
也是不变子空间，因为置换不改变坐标和。并且
\[
\C^3=U\oplus W.
\]
这说明该置换表示是可约的。
\end{example}

我们希望知道：给定一个子表示 $W\subseteq V$，是否总能找到另一个子表示 $U$ 使
\[
V=W\oplus U?
\]
在线性代数中，任意子空间总有线性补空间；但这个补空间未必在群作用下保持不变。Maschke 定理说明：当系数域特征不整除群的阶时，可以把任意补空间“平均化”为不变补空间。

\begin{proposition}
设 $W\subseteq V$ 是表示 $V$ 的一个子表示。则下列两条件等价：
\begin{enumerate}[label=\textnormal{(\arabic*)}]
  \item 存在子表示 $U\subseteq V$ 使得 $V=W\oplus U$；
  \item 存在线性映射 $P:V\to V$，满足 $P^2=P$，$\Ima(P)=W$，且 $P$ 与群作用可交换，即
  \[
  P\rho(g)=\rho(g)P\qquad (g\in G).
  \]
\end{enumerate}
\end{proposition}

\begin{proof}
若存在 $V=W\oplus U$ 且 $U$ 为子表示，则定义投影 $P$ 为沿 $U$ 到 $W$ 的投影。显然 $P^2=P$ 且像为 $W$。对任意 $v=w+u$（$w\in W,u\in U$）及 $g\in G$，有
\[
P\rho(g)(v)=P(\rho(g)w+\rho(g)u)=\rho(g)w,
\]
因为 $W,U$ 都不变；另一方面
\[
\rho(g)P(v)=\rho(g)w.
\]
故 $P\rho(g)=\rho(g)P$。

反过来，若存在满足条件的 $P$，令 $U=\Ker(P)$。因为 $P^2=P$，线性代数告诉我们
\[
V=\Ima(P)\oplus \Ker(P)=W\oplus U.
\]
只需证明 $U$ 为子表示。若 $u\in U$，则 $P(u)=0$。对任意 $g\in G$，有
\[
P(\rho(g)u)=\rho(g)P(u)=0,
\]
故 $\rho(g)u\in \Ker(P)=U$。因此 $U$ 也是子表示。
\end{proof}

\begin{theorem}[Maschke 定理]
设 $G$ 为有限群，$F$ 为域，且 $\Char(F)\nmid |G|$。则 $G$ 的任意有限维表示都是完全可约的：换言之，任意子表示都有不变补空间。
\end{theorem}

\begin{proof}
设 $\rho:G\to \GL(V)$ 是一个表示，$W\subseteq V$ 是任意子表示。先任取一个线性补空间 $V=W\oplus U_0$；注意 $U_0$ 不一定是子表示。令
\[
P_0:V\to V
\]
为沿 $U_0$ 到 $W$ 的投影，则 $P_0^2=P_0$，且 $\Ima(P_0)=W$。

我们把 $P_0$ 在群上平均：定义
\[
P=\frac{1}{|G|}\sum_{g\in G}\rho(g)P_0\rho(g)^{-1}.
\]
由于 $\Char(F)\nmid |G|$，元素 $|G|$ 在 $F$ 中可逆，因此上式有意义。

先证 $P$ 与群作用可交换。对任意 $h\in G$，
\[
\rho(h)P\rho(h)^{-1}
=\frac{1}{|G|}\sum_{g\in G}\rho(hg)P_0\rho(hg)^{-1}.
\]
当 $g$ 遍历 $G$ 时，$hg$ 也遍历 $G$，故上式等于 $P$。于是
\[
\rho(h)P=P\rho(h).
\]

再证 $\Ima(P)=W$ 且 $P$ 在 $W$ 上恒等。若 $w\in W$，因 $W$ 为子表示，$\rho(g)^{-1}w\in W$，而 $P_0$ 在 $W$ 上等于恒等，故
\[
\rho(g)P_0\rho(g)^{-1}(w)=\rho(g)\rho(g)^{-1}(w)=w.
\]
对所有 $g$ 求和得 $P(w)=w$。因此 $P$ 在 $W$ 上恒等，特别地 $W\subseteq \Ima(P)$。

另一方面，对任意 $v\in V$，每个向量 $\rho(g)P_0\rho(g)^{-1}(v)$ 都属于 $W$，因为 $P_0$ 的像是 $W$，而 $W$ 又在群作用下不变。故它们的平均 $P(v)$ 仍在 $W$ 中，从而 $\Ima(P)\subseteq W$。

于是 $\Ima(P)=W$，且 $P^2=P$（因为 $P$ 在其像 $W$ 上为恒等）。由上一命题，$W$ 有不变补空间，定理得证。
\end{proof}

\begin{corollary}
在 Maschke 定理的假设下，任意有限维表示都可以分解为有限个不可约表示的直和。
\end{corollary}

\begin{proof}
对 $\dim V$ 作归纳。若 $V$ 不可约，则已完成。若 $V$ 可约，则存在非平凡子表示 $0\neq W\neq V$。由 Maschke 定理，存在子表示 $U$ 使 $V=W\oplus U$。又有
\[
\dim W<\dim V,\qquad \dim U<\dim V.
\]
归纳假设说明 $W,U$ 都可分解为不可约表示直和，因此 $V$ 也可分解为不可约表示直和。
\end{proof}

\begin{remark}
上推论中的分解通常不是唯一的；但在复数域上，不可约表示出现的\emph{重数}是唯一的。这个唯一性将在特征标理论中自然体现出来。
\end{remark}

\begin{example}[Maschke 定理失效的情形]
设 $G=C_p=\langle g\mid g^p=1\rangle$，系数域取 $F=\F_p$。令 $V=F^2$，并定义
\[
\rho(g)=\mat{1&1\\0&1}.
\]
因为在特征 $p$ 中有
\[
\rho(g)^p=\mat{1&1\\0&1}^p=I_2,
\]
故这确实给出 $C_p$ 的一个表示。设
\[
W=F\mat{1\\0}.
\]
则 $W$ 是子表示，但它没有不变补空间。事实上，若 $U=F\mat{a\\1}$ 是任一与 $W$ 互补的一维子空间，则
\[
\rho(g)\mat{a\\1}=\mat{a+1\\1},
\]
这不在 $U$ 中，因为它不是 $\mat{a\\1}$ 的标量倍。故当 $\Char(F)\mid |G|$ 时，完全可约性可能失败。
\end{example}

\begin{proposition}
设 $F=\C$。则任意有限群表示 $\rho:G\to \GL(V)$ 都可以在某个 Hermite 内积下化为酉表示，即存在正定 Hermite 内积 $\inner{\cdot}{\cdot}_G$ 使得
\[
\inner{\rho(g)v}{\rho(g)w}_G=\inner{v}{w}_G
\qquad (g\in G,\,v,w\in V).
\]
\end{proposition}

\begin{proof}
先任取 $V$ 上一个 Hermite 内积 $\inner{\cdot}{\cdot}_0$。定义新的内积
\[
\inner{v}{w}_G=\frac{1}{|G|}\sum_{g\in G}\inner{\rho(g)v}{\rho(g)w}_0.
\]
这是有限个 Hermite 内积的平均，因此仍是 Hermite 内积。对任意 $h\in G$，
\[
\inner{\rho(h)v}{\rho(h)w}_G
=\frac{1}{|G|}\sum_{g\in G}\inner{\rho(gh)v}{\rho(gh)w}_0.
\]
因为当 $g$ 遍历 $G$ 时，$gh$ 也遍历 $G$，所以上式恰等于 $\inner{v}{w}_G$。故群作用在此内积下为酉变换。
\end{proof}

\section{Schur 引理}

Schur 引理是不可约表示论的核心工具。它说明：不可约表示之间的 $G$-同态极少；在同一个不可约表示内部，可与群作用可交换的线性变换也极少。

\begin{theorem}[Schur 引理]
设 $V,W$ 是域 $F$ 上的两个不可约表示，$T:V\to W$ 是一个 $G$-同态，即
\[
T\rho_V(g)=\rho_W(g)T\qquad (g\in G).
\]
则：
\begin{enumerate}[label=\textnormal{(\arabic*)}]
  \item 若 $T\neq 0$，则 $T$ 是同构；
  \item 因而若 $V\not\cong W$，则 $\Hom_G(V,W)=0$；
  \item 当 $V=W$ 时，$\End_G(V)$ 是一个除环。
\end{enumerate}
\end{theorem}

\begin{proof}
因为 $T$ 与群作用可交换，所以
\[
\Ker(T)\subseteq V,\qquad \Ima(T)\subseteq W
\]
都是不变子空间。若 $T\neq 0$，则 $\Ima(T)\neq 0$。由于 $W$ 不可约，只能有 $\Ima(T)=W$，故 $T$ 满射。

另一方面，$\Ker(T)$ 是 $V$ 的不变子空间。若 $\Ker(T)=V$，则 $T=0$，矛盾；因此只能有 $\Ker(T)=0$。于是 $T$ 单射。

综上，$T$ 为双射，即同构。

第二条立即推出：若 $V,W$ 不等价，则不存在非零同态。

最后设 $V=W$，取任意非零 $T\in \End_G(V)$。由第一条，$T$ 必为同构，因此在 $\End_G(V)$ 中可逆。故 $\End_G(V)$ 的每个非零元都可逆，它是一个除环。
\end{proof}

\begin{corollary}
设 $V$ 是不可约表示，若线性变换 $T\in \End(V)$ 与所有 $\rho(g)$ 都可交换，则 $T$ 要么为零，要么为同构。
\end{corollary}

\begin{proof}
这只是 Schur 引理在 $W=V$ 时的直接表述。
\end{proof}

\begin{corollary}
设 $F=\C$，$V$ 是一个有限维不可约复表示，则
\[
\End_G(V)=\C\Id_V.
\]
换言之，任何与群作用可交换的线性变换都只是数乘。
\end{corollary}

\begin{proof}
取任意 $T\in \End_G(V)$。复数域代数闭，因此 $T$ 至少有一个特征值 $\lambda\in \C$。于是
\[
T-\lambda \Id_V
\]
不是同构。又它仍与群作用可交换，所以属于 $\End_G(V)$。由上一个推论，$T-\lambda\Id_V$ 只能为零。故 $T=\lambda\Id_V$。
\end{proof}

\begin{proposition}
设 $V_1,\dots,V_r$ 是两两不等价的不可约表示，$m_1,\dots,m_r\in \N$。令
\[
V\cong \bigoplus_{i=1}^r V_i^{\oplus m_i}.
\]
则
\[
\End_G(V)\cong \bigoplus_{i=1}^r M_{m_i}(\End_G(V_i)).
\]
特别地，在 $F=\C$ 时，
\[
\dim_{\C}\End_G(V)=\sum_{i=1}^r m_i^2.
\]
\end{proposition}

\begin{proof}
不同 $V_i$ 之间没有非零 $G$-同态，因此任一 $G$-自同态必须保持每个等价类对应的直和块。于是
\[
\End_G(V)\cong \bigoplus_{i=1}^r \End_G(V_i^{\oplus m_i}).
\]
而 $V_i^{\oplus m_i}$ 的 $G$-自同态正是 $m_i\times m_i$ 矩阵，其条目属于 $\End_G(V_i)$。故得上式。

在复数域上，Schur 引理给出 $\End_G(V_i)=\C\Id$，故第 $i$ 块同构于 $M_{m_i}(\C)$，维数为 $m_i^2$。总维数即为各块维数之和。
\end{proof}

\begin{remark}
Schur 引理在后面会反复出现。它的一个典型用法是：只要我们构造出一个与群作用可交换的算子，便能借助不可约性立刻把它压缩为一个标量。这正是特征标正交关系证明中的关键步骤。
\end{remark}

\section{特征标理论初步}

从现在开始，我们取系数域为 $\C$。由于复表示可以酉化，我们可以放心地使用迹、共轭和内积等工具。

\begin{definition}
设 $\rho:G\to \GL(V)$ 是一个有限维复表示。定义其\emph{特征标}（character）为函数
\[
\chi_\rho:G\to \C,\qquad \chi_\rho(g)=\Tr(\rho(g)).
\]
若上下文明确，也简记为 $\chi$。
\end{definition}

\begin{proposition}
特征标是类函数，即若 $g,h\in G$ 共轭，则
\[
\chi(hgh^{-1})=\chi(g).
\]
\end{proposition}

\begin{proof}
由表示同态性质，
\[
\rho(hgh^{-1})=\rho(h)\rho(g)\rho(h)^{-1}.
\]
相似矩阵有相同的迹，因此
\[
\chi(hgh^{-1})=\Tr(\rho(h)\rho(g)\rho(h)^{-1})=\Tr(\rho(g))=\chi(g).
\]
\end{proof}

\begin{proposition}
若 $\rho=\rho_1\oplus \rho_2$，则
\[
\chi_\rho=\chi_{\rho_1}+\chi_{\rho_2}.
\]
若 $V,W$ 是表示，则张量积表示 $V\otimes W$ 的特征标满足
\[
\chi_{V\otimes W}(g)=\chi_V(g)\chi_W(g).
\]
\end{proposition}

\begin{proof}
对直和而言，$\rho(g)$ 在适当基下是分块对角矩阵，故迹等于各块迹之和。

对张量积而言，若 $\rho_V(g)$ 与 $\rho_W(g)$ 的特征值分别为 $\lambda_1,\dots,\lambda_m$ 与 $\mu_1,\dots,\mu_n$，则 $\rho_{V\otimes W}(g)=\rho_V(g)\otimes \rho_W(g)$ 的特征值为所有 $\lambda_i\mu_j$。于是
\[
\chi_{V\otimes W}(g)=\sum_{i,j}\lambda_i\mu_j
=\left(\sum_i\lambda_i\right)\left(\sum_j\mu_j\right)
=\chi_V(g)\chi_W(g).
\]
\end{proof}

\begin{definition}
若 $f_1,f_2:G\to \C$ 是两个类函数，定义它们的内积为
\[
\inner{f_1}{f_2}=\frac{1}{|G|}\sum_{g\in G}f_1(g)\overline{f_2(g)}.
\]
\end{definition}

这个内积之所以自然，是因为它与正交关系、重数公式和分解唯一性完全兼容。

为了证明正交关系，我们先建立矩阵系数的正交性。设 $V,W$ 是不可约酉表示，分别取标准正交基 $e_1,\dots,e_m$ 与 $f_1,\dots,f_n$。记
\[
\rho_{ab}(g)=\inner{\rho(g)e_b}{e_a},\qquad
\sigma_{cd}(g)=\inner{\sigma(g)f_d}{f_c}.
\]

\begin{lemma}[矩阵系数正交关系]
在上述记号下，有
\[
\sum_{g\in G}\sigma_{cj}(g)\overline{\rho_{bi}(g)}=
\begin{cases}
0, & V\not\cong W,\\[4pt]
\dfrac{|G|}{\dim V}\,\delta_{cb}\delta_{ji}, & V\cong W.
\end{cases}
\]
\end{lemma}

\begin{proof}
对任意线性映射 $T:V\to W$，定义
\[
\Phi(T)=\sum_{g\in G}\sigma(g)T\rho(g)^{-1}.
\]
对任意 $h\in G$，
\[
\sigma(h)\Phi(T)\rho(h)^{-1}
=\sum_{g\in G}\sigma(hg)T\rho(hg)^{-1}=
\Phi(T),
\]
故 $\Phi(T)\in \Hom_G(V,W)$。

若 $V\not\cong W$，由 Schur 引理得 $\Phi(T)=0$。若 $V\cong W$，特别地可把 $V,W$ 视为同一不可约表示，则 $\Phi(T)$ 与群作用可交换，因此由 Schur 引理可知
\[
\Phi(T)=\lambda(T)\Id_V.
\]
对两边取迹，得到
\[
\dim V\cdot \lambda(T)=\Tr(\Phi(T))
=\sum_{g\in G}\Tr(\rho(g)T\rho(g)^{-1})
=\sum_{g\in G}\Tr(T)=|G|\Tr(T).
\]
故
\[
\Phi(T)=\frac{|G|\Tr(T)}{\dim V}\Id_V.
\]

现在取 $T=T_{ji}:V\to W$ 为线性映射：$T_{ji}(e_i)=f_j$，其余基向量送到 $0$。则对任意指标 $b,c$，
\[
\inner{\Phi(T_{ji})e_b}{f_c}
=\sum_{g\in G}\inner{\sigma(g)T_{ji}\rho(g)^{-1}e_b}{f_c}
=\sum_{g\in G}\sigma_{cj}(g)\overline{\rho_{bi}(g)}.
\]
若 $V\not\cong W$，左边为 $0$。若 $V\cong W$，则 $\Tr(T_{ji})=\delta_{ij}$，所以
\[
\Phi(T_{ji})=\frac{|G|\delta_{ij}}{\dim V}\Id_V,
\]
从而
\[
\inner{\Phi(T_{ji})e_b}{f_c}=\frac{|G|}{\dim V}\delta_{ij}\delta_{bc}.
\]
于是结论成立。
\end{proof}

\begin{theorem}[第一正交关系]
设 $\chi,\psi$ 是两个不可约复表示的特征标，则
\[
\inner{\chi}{\psi}=\frac{1}{|G|}\sum_{g\in G}\chi(g)\overline{\psi(g)}=
\begin{cases}
1,& \chi=\psi,\\
0,& \chi\neq \psi.
\end{cases}
\]
换言之，不同不可约特征标彼此正交，且每个不可约特征标的范数为 $1$。
\end{theorem}

\begin{proof}
设 $\chi,\psi$ 分别来自不可约表示 $V,W$，其维数分别为 $m,n$。将特征标写成矩阵对角元之和：
\[
\chi(g)=\sum_{a=1}^m \rho_{aa}(g),
\qquad
\psi(g)=\sum_{b=1}^n \sigma_{bb}(g).
\]
于是
\[
\sum_{g\in G}\chi(g)\overline{\psi(g)}
=\sum_{a,b}\sum_{g\in G}\rho_{aa}(g)\overline{\sigma_{bb}(g)}.
\]
若 $V\not\cong W$，由上一个引理，每一项都为 $0$，故总和为 $0$。

若 $V\cong W$，则 $m=n$，并且
\[
\sum_{g\in G}\rho_{aa}(g)\overline{\rho_{bb}(g)}
=\frac{|G|}{m}\delta_{ab}.
\]
把 $a,b$ 求和得
\[
\sum_{g\in G}|\chi(g)|^2
=\sum_{a,b}\frac{|G|}{m}\delta_{ab}=|G|.
\]
两边除以 $|G|$ 即得结论。
\end{proof}

\begin{proposition}[正则表示的特征标]
设 $\chi_{\mathrm{reg}}$ 为正则表示的特征标，则
\[
\chi_{\mathrm{reg}}(g)=
\begin{cases}
|G|, & g=1,\\
0, & g\neq 1.
\end{cases}
\]

并且若 $V_i$ 的不可约特征标为 $\chi_i$、维数为 $d_i$，则在正则表示分解中，$V_i$ 出现的重数恰为 $d_i$。
\end{proposition}

\begin{proof}
正则表示下，$g$ 对基 $\{e_x\}$ 的作用为 $e_x\mapsto e_{gx}$。若 $g\neq 1$，则左乘 $x\mapsto gx$ 没有不动点，因此矩阵对角元全为 $0$，故迹为 $0$。若 $g=1$，则矩阵为恒等矩阵，迹为 $|G|$。

设正则表示分解为
\[
\C[G]\cong \bigoplus_i m_iV_i.
\]
其特征标满足
\[
\chi_{\mathrm{reg}}=\sum_i m_i\chi_i.
\]
对两边与 $\chi_j$ 取内积，利用第一正交关系得
\[
m_j=\inner{\chi_{\mathrm{reg}}}{\chi_j}
=\frac{1}{|G|}\sum_{g\in G}\chi_{\mathrm{reg}}(g)\overline{\chi_j(g)}
=\overline{\chi_j(1)}=\chi_j(1)=d_j.
\]
这里用到了 $\chi_j(1)=\dim V_j$ 为正整数。
\end{proof}

由上命题立刻得到一个重要恒等式：
\[
|G|=\dim \C[G]=\sum_i d_i\cdot m_i=\sum_i d_i^2.
\]

接下来说明不可约特征标事实上张成了全部类函数空间。

\begin{theorem}
不可约特征标构成类函数空间的一组正交规范基。特别地，不可约表示的个数等于 $G$ 的共轭类个数。
\end{theorem}

\begin{proof}
第一正交关系已经说明不可约特征标两两正交，因此它们线性无关。只需证明它们张成全部类函数。

设 $V_1,\dots,V_r$ 为两两不等价的不可约表示，$d_i=\dim V_i$。由上一命题，正则表示中 $V_i$ 出现 $d_i$ 次，从而
\[
|G|=\sum_{i=1}^r d_i^2.
\]
另一方面，矩阵系数正交关系说明：对每个 $i$ 取一组标准正交基后，所有函数
\[
\rho^{(i)}_{ab}:G\to \C \qquad (1\le a,b\le d_i,
\ 1\le i\le r)
\]
两两正交，且总数是 $\sum_i d_i^2=|G|$。由于 $G$ 上所有复值函数构成的向量空间维数正是 $|G|$，故这些矩阵系数构成该函数空间的一组基。

现在设 $f:G\to \C$ 为任意类函数，把它写成矩阵系数线性组合。只需说明每个矩阵系数在按共轭平均后只留下特征标成分。固定某个不可约表示 $\rho$ 及矩阵单位 $E_{ba}$，注意
\[
\rho_{ab}(x)=\Tr(E_{ba}\rho(x)).
\]
考虑其类平均
\[
\rho_{ab}^{\#}(x)=\frac{1}{|G|}\sum_{t\in G}\rho_{ab}(txt^{-1}).
\]
将其写成迹的形式：
\[
\rho_{ab}^{\#}(x)=\Tr\left(\frac{1}{|G|}\sum_{t\in G}\rho(t)^{-1}E_{ba}\rho(t)\cdot \rho(x)\right).
\]
令
\[
A=\frac{1}{|G|}\sum_{t\in G}\rho(t)^{-1}E_{ba}\rho(t).
\]
则对任意 $h\in G$，$\rho(h)A\rho(h)^{-1}=A$，所以 $A\in \End_G(V)$。由 Schur 引理，$A=\lambda I$。取迹得
\[
\lambda d=\Tr(A)=\Tr(E_{ba})=\delta_{ab},
\]
故
\[
A=\frac{\delta_{ab}}{d}I.
\]
于是
\[
\rho_{ab}^{\#}(x)=\frac{\delta_{ab}}{d}\chi_\rho(x).
\]
因此任何类函数都可以写成不可约特征标的线性组合。

既然不可约特征标既正交又张成类函数空间，它们就是一组正交规范基。类函数空间维数等于共轭类个数，故不可约特征标的个数也等于共轭类个数。
\end{proof}

\begin{theorem}[第二正交关系]
设 $g_1,\dots,g_r$ 是 $G$ 的共轭类代表元，$C_i$ 为 $g_i$ 的共轭类，$Z_i=C_G(g_i)$ 为中心化子。则对任意 $1\le i,j\le r$，
\[
\sum_{\chi\ \mathrm{irr}} \chi(g_i)\overline{\chi(g_j)}=
\begin{cases}
|Z_i|,& i=j,\\
0,& i\neq j.
\end{cases}
\]
其中求和遍历全部不可约特征标。
\end{theorem}

\begin{proof}
令 $\delta_i$ 为共轭类 $C_i$ 的示性函数：
\[
\delta_i(g)=
\begin{cases}
1,& g\in C_i,\\
0,& g\notin C_i.
\end{cases}
\]
它是一个类函数。由于不可约特征标构成类函数空间的正交规范基，故
\[
\delta_i=\sum_{\chi\ \mathrm{irr}} \inner{\delta_i}{\chi}\,\chi.
\]
而
\[
\inner{\delta_i}{\chi}
=\frac{1}{|G|}\sum_{g\in G}\delta_i(g)\overline{\chi(g)}
=\frac{|C_i|}{|G|}\overline{\chi(g_i)}.
\]
因此
\[
\delta_i(g_j)=\frac{|C_i|}{|G|}\sum_{\chi\ \mathrm{irr}}\overline{\chi(g_i)}\chi(g_j).
\]
若 $i\neq j$，左边为 $0$；若 $i=j$，左边为 $1$。故
\[
\sum_{\chi\ \mathrm{irr}}\chi(g_i)\overline{\chi(g_j)}=
\begin{cases}
\dfrac{|G|}{|C_i|},& i=j,\\
0,& i\neq j.
\end{cases}
\]
而 $|G|=|C_i|\cdot |C_G(g_i)|$，所以 $|G|/|C_i|=|Z_i|$。结论得证。
\end{proof}

\begin{remark}
第一正交关系说“特征标的行彼此正交”，第二正交关系说“特征标表的列彼此正交”。这是特征标表之所以强大的根本原因。
\end{remark}

\begin{example}[循环群 $\Z/n\Z$ 的特征标]
设 $G=\langle t\mid t^n=1\rangle\cong \Z/n\Z$，取一个本原 $n$ 次单位根 $\zeta_n\in \C$。对每个 $j=0,1,\dots,n-1$，定义
\[
\chi_j(t^m)=\zeta_n^{jm}.
\]
这是一个一维表示的特征标，因为
\[
\chi_j(t^{m_1+m_2})=\zeta_n^{j(m_1+m_2)}=
\chi_j(t^{m_1})\chi_j(t^{m_2}).
\]
反过来，$G$ 的任意不可约复表示都是一维的：因为 $G$ 阿贝尔，所以对任意 $g\in G$，$\rho(g)$ 与所有 $\rho(h)$ 可交换；若 $\rho$ 不可约，Schur 引理推出每个 $\rho(g)$ 都是标量，于是表示空间只能是一维。故全部不可约特征标恰为 $\chi_0,\dots,\chi_{n-1}$。

它们满足
\[
\frac{1}{n}\sum_{m=0}^{n-1}\chi_i(t^m)\overline{\chi_j(t^m)}=\delta_{ij}.
\]
这正是有限 Fourier 分析的正交关系。
\end{example}

\begin{example}[$S_3$ 的特征标表]
群 $S_3$ 有三个共轭类：
\[
\{1\},\qquad \{(12),(13),(23)\},\qquad \{(123),(132)\}.
\]
因此有三个不可约表示。已知两个一维表示：平凡表示 $\mathbf{1}$ 与符号表示 $\sgn$。第三个表示来自三维置换表示去掉平凡子表示后的二维部分，记为 $\mathrm{Std}$（标准表示）。

三维置换表示的特征标由“不动点个数”给出，因此在三个共轭类上的值分别为 $3,1,0$。去掉平凡表示后，$\mathrm{Std}$ 的特征标为
\[
\chi_{\mathrm{Std}}=(2,0,-1).
\]
于是 $S_3$ 的特征标表为
\[
\begin{array}{c|ccc}
 & 1 & (12) & (123)\\
\hline
\mathbf{1} & 1 & 1 & 1\\
\sgn & 1 & -1 & 1\\
\mathrm{Std} & 2 & 0 & -1
\end{array}
\]
可直接检查第一正交关系：
\[
\frac{1}{6}(1\cdot 1+3\cdot 1\cdot (-1)+2\cdot 1\cdot 1)=0,
\]
等等。
\end{example}

\begin{example}[$D_4$ 的特征标表]
设
\[
D_4=\langle r,s\mid r^4=s^2=1,\ srs=r^{-1}\rangle
\]
为正方形的对称群，$|D_4|=8$。它的共轭类为
\[
\{1\},\ \{r^2\},\ \{r,r^3\},\ \{s,sr^2\},\ \{sr,sr^3\}.
\]
因此有五个不可约表示。

首先，$D_4$ 的交换化同构于 $C_2\times C_2$，所以有四个一维表示。它们由 $r,s$ 的像决定，且都满足 $r^2\mapsto 1$。记这四个特征标分别为
\[
\chi_{\varepsilon,\delta}(r)=\varepsilon,\qquad \chi_{\varepsilon,\delta}(s)=\delta,
\qquad \varepsilon,\delta\in\{\pm 1\}.
\]
此外，正方形在平面上的几何作用给出一个二维表示 $\rho_{\mathrm{geo}}$：
\[
\rho_{\mathrm{geo}}(r)=\mat{0&-1\\1&0},
\qquad
\rho_{\mathrm{geo}}(s)=\mat{1&0\\0&-1}.
\]
其特征标在五个共轭类上的值分别为
\[
2,\ -2,\ 0,\ 0,\ 0.
\]
由平方和公式
\[
1^2+1^2+1^2+1^2+2^2=8
\]
可知这已经给出全部不可约表示。特征标表如下：
\[
\begin{array}{c|ccccc}
 & 1 & r^2 & r,r^3 & s,sr^2 & sr,sr^3\\
\hline
\chi_{1,1} & 1 & 1 & 1 & 1 & 1\\
\chi_{1,-1} & 1 & 1 & 1 & -1 & -1\\
\chi_{-1,1} & 1 & 1 & -1 & 1 & -1\\
\chi_{-1,-1} & 1 & 1 & -1 & -1 & 1\\
\rho_{\mathrm{geo}} & 2 & -2 & 0 & 0 & 0
\end{array}
\]
\end{example}

\section{$\GL_2$ 与 $\SL_2$ 的表示}

本节给出两个方向的例子。第一类是代数表示：由 $\GL_2$ 或 $\SL_2$ 对多项式空间的自然作用得到对称幂表示 $\mathrm{Sym}^k$。第二类是有限域上 $\GL_2(\F_q)$ 的有限群表示：其中主系列与 Steinberg 表示是最基本的对象。

\subsection*{对称幂表示}

设 $F$ 为任意域，$V=F^2$ 是 $\GL_2(F)$ 的自然表示。对每个整数 $k\ge 0$，记 $\mathrm{Sym}^k(V)$ 为 $V$ 的第 $k$ 个对称幂；更具体地，它可以与所有 $X,Y$ 的 $k$ 次齐次多项式空间同一：
\[
\mathrm{Sym}^k(V)\cong F[X,Y]_k=
\left\{\sum_{i=0}^k a_iX^{k-i}Y^i\right\}.
\]

\begin{proposition}
对任意 $g=\mat{a&b\\c&d}\in \GL_2(F)$，定义
\[
(g\cdot P)(X,Y)=P(aX+cY,bX+dY),\qquad P\in F[X,Y]_k.
\]
则这给出 $\GL_2(F)$ 在 $F[X,Y]_k$ 上的一个 $(k+1)$ 维表示，称为\emph{第 $k$ 个对称幂表示}，记作 $\mathrm{Sym}^k$。
\end{proposition}

\begin{proof}
首先，若 $P$ 是 $k$ 次齐次多项式，则把 $X,Y$ 换成线性组合 $aX+cY$ 与 $bX+dY$ 后，所得仍是 $k$ 次齐次多项式，所以作用确实落在 $F[X,Y]_k$ 内。

再验证表示律。若
\[
g=\mat{a&b\\c&d},\qquad h=\mat{a'&b'\\c'&d'},
\]
则
\[
(h\cdot P)(X,Y)=P(a'X+c'Y,b'X+d'Y).
\]
接着作用 $g$，得
\[
(g\cdot (h\cdot P))(X,Y)
=P\bigl(a'(aX+cY)+c'(bX+dY),\ b'(aX+cY)+d'(bX+dY)\bigr).
\]
括号中的线性替换恰对应矩阵乘积 $gh$。因我们约定右边代入的是行向量 $(X,Y)$ 乘矩阵，所以这正等于 $(gh)\cdot P$。此外 $I_2$ 显然作用为恒等。故这是一表示。

空间 $F[X,Y]_k$ 的标准基可取为
\[
X^k,X^{k-1}Y,\dots,Y^k,
\]
故其维数为 $k+1$。
\end{proof}

\begin{example}
$k=0$ 时，$\mathrm{Sym}^0$ 是平凡表示；
$k=1$ 时，$\mathrm{Sym}^1$ 就是自然表示；
$k=2$ 时，表示空间基可取 $X^2,XY,Y^2$，而
\[
\mat{a&b\\c&d}\cdot X^2=(aX+cY)^2=a^2X^2+2ac\,XY+c^2Y^2,
\]
因此矩阵可显式写出为
\[
\mat{a^2&ab&b^2\\2ac&ad+bc&2bd\\c^2&cd&d^2}
\]
（基的排列依约定略作调整即可）。
\end{example}

\begin{remark}
将 $\mathrm{Sym}^k$ 限制到 $\SL_2(F)$ 上，仍得到 $\SL_2(F)$ 的表示。对模形式来说，权 $k+2$ 往往与 $\mathrm{Sym}^k$ 对应；这也是 $\GL_2$ 表示不断出现于模形式理论中的原因之一。
\end{remark}

\subsection*{$\GL_2(\F_q)$ 的主系列与 Steinberg 表示}

从现在起令 $q$ 为素数幂，
\[
G=\GL_2(\F_q),
\qquad
B=\left\{\mat{a&b\\0&d}:a,d\in \F_q^{\times},\ b\in \F_q\right\}
\]
为上三角 Borel 子群。

记 $\mathbf{P}^1(\F_q)$ 为 $\F_q^2$ 中一维子空间的集合。$G$ 自然作用在 $\mathbf{P}^1(\F_q)$ 上，且作用是传递的。基点 $[1:0]$ 的稳定子正是 $B$，故有集合双射
\[
G/B\cong \mathbf{P}^1(\F_q).
\]

\begin{proposition}
置换表示 $\C[\mathbf{P}^1(\F_q)]$ 同构于从 $B$ 的平凡表示诱导上来的表示
\[
\mathrm{Ind}_B^G \mathbf{1}.
\]
其维数为 $q+1$。
\end{proposition}

\begin{proof}
按诱导表示的定义，$\mathrm{Ind}_B^G\mathbf{1}$ 可以看成满足
\[
 f(bg)=f(g)\qquad (b\in B,g\in G)
\]
的函数 $f:G\to \C$ 的空间。这样的函数完全由右陪集 $B\backslash G$ 上的取值决定；而取逆映射给出 $B\backslash G\cong G/B$，又 $G/B\cong \mathbf{P}^1(\F_q)$，因此该空间同构于 $\mathbf{P}^1(\F_q)$ 上的复值函数空间，也就是置换表示 $\C[\mathbf{P}^1(\F_q)]$。

因为 $|\mathbf{P}^1(\F_q)|=q+1$，所以维数为 $q+1$。
\end{proof}

\begin{definition}
定义
\[
\mathrm{St}=
\left\{f:\mathbf{P}^1(\F_q)\to \C\ :\ \sum_{x\in \mathbf{P}^1(\F_q)}f(x)=0\right\}.
\]
这是一个维数为 $q$ 的 $G$-不变子空间，称为 $\GL_2(\F_q)$ 的\emph{Steinberg 表示}。
\end{definition}

\begin{proposition}
有直和分解
\[
\C[\mathbf{P}^1(\F_q)]\cong \mathbf{1}\oplus \mathrm{St}.
\]
其中 $\mathbf{1}$ 是常值函数构成的一维平凡子表示。并且 $\mathrm{St}$ 是不可约表示。
\end{proposition}

\begin{proof}
常值函数显然构成一维平凡子表示。任意函数 $f$ 都可唯一写成
\[
f=\left(f-\frac{1}{q+1}\sum_x f(x)\right)+\frac{1}{q+1}\sum_x f(x),
\]
其中第一项属于 $\mathrm{St}$，第二项为常值函数，因此
\[
\C[\mathbf{P}^1(\F_q)]=\mathbf{1}\oplus \mathrm{St}.
\]

接下来证 $\mathrm{St}$ 不可约。考虑任意 $G$-自同态
\[
A:\C[\mathbf{P}^1(\F_q)]\to \C[\mathbf{P}^1(\F_q)].
\]
记其在点基 $\{e_x:x\in \mathbf{P}^1(\F_q)\}$ 下的矩阵系数为 $a_{xy}$。由于 $A$ 与 $G$ 作用可交换，而 $G$ 在有序点对集合上有两条轨道：
\[
\{(x,x)	ext{ 的全体}\},\qquad \{(x,y):x\neq y\},
\]
故存在常数 $\alpha,\beta\in \C$ 使得
\[
a_{xx}=\alpha,\qquad a_{xy}=\beta\quad (x\neq y).
\]
因此
\[
A=(\alpha-\beta)I+\beta J,
\]
其中 $J$ 是所有矩阵元都为 $1$ 的矩阵。

现在把 $A$ 限制到 $\mathrm{St}$ 上。对 $f\in \mathrm{St}$，有各坐标和为 $0$，所以 $Jf=0$。于是
\[
A|_{\mathrm{St}}=(\alpha-\beta)\Id_{\mathrm{St}}.
\]
这说明
\[
\End_G(\mathrm{St})=\C\Id.
\]

若 $\mathrm{St}$ 可约，则按 Maschke 定理可分解为两个非零子表示的直和。沿某一块的投影将给出一个既非零也非标量的 $G$-自同态，矛盾。故 $\mathrm{St}$ 不可约。
\end{proof}

\begin{definition}
设 $\chi_1,\chi_2:\F_q^{\times}\to \C^{\times}$ 为两个乘法特征。定义 $B$ 的一维表示
\[
\chi_1\boxtimes \chi_2\left(\mat{a&b\\0&d}\right)=\chi_1(a)\chi_2(d).
\]
从它诱导得到的表示
\[
I(\chi_1,\chi_2)=\mathrm{Ind}_B^G(\chi_1\boxtimes \chi_2)
\]
称为 $\GL_2(\F_q)$ 的一个\emph{主系列表示}（principal series）。
\end{definition}

\begin{proposition}
对任意乘法特征 $\chi_1,\chi_2$，主系列表示 $I(\chi_1,\chi_2)$ 的维数都是 $q+1$。
\end{proposition}

\begin{proof}
与平凡诱导表示的情形相同，$I(\chi_1,\chi_2)$ 可看成满足
\[
 f(bg)=(\chi_1\boxtimes\chi_2)(b)f(g)
\]
的函数空间。这样的函数由 $B\backslash G$ 上的取值唯一决定，所以维数等于陪集数 $[G:B]=|\mathbf{P}^1(\F_q)|=q+1$。
\end{proof}

\begin{proposition}
若 $\chi:\F_q^{\times}\to \C^{\times}$ 是乘法特征，则
\[
I(\chi,\chi)\cong (\chi\circ \det)\otimes I(1,1)
\cong (\chi\circ \det)\oplus \bigl((\chi\circ \det)\otimes \mathrm{St}\bigr).
\]
后一个 $q$ 维表示通常记作 $\mathrm{St}_\chi$，称为扭曲的 Steinberg 表示。
\end{proposition}

\begin{proof}
对满足 $f(bg)=\chi(a)\chi(d)f(g)$ 的函数 $f$，定义
\[
\Phi(f)(g)=\chi(\det g)^{-1}f(g).
\]
因为
\[
\det(bg)=\det(b)\det(g)=ad\det(g),
\]
故
\[
\Phi(f)(bg)=\chi(ad)^{-1}\chi(a)\chi(d)f(g)=\Phi(f)(g).
\]
所以 $\Phi$ 给出 $I(\chi,\chi)\to I(1,1)$ 的线性同构。检查 $G$ 作用可知，它正对应于张量上 $\chi\circ\det$。于是
\[
I(\chi,\chi)\cong (\chi\circ\det)\otimes I(1,1).
\]
再利用前面证明的分解 $I(1,1)=\mathbf{1}\oplus \mathrm{St}$，对两边张量 $\chi\circ\det$ 即得
\[
I(\chi,\chi)\cong (\chi\circ\det)\oplus ((\chi\circ\det)\otimes \mathrm{St}).
\]
\end{proof}

\begin{remark}
当 $\chi_1\neq \chi_2$ 时，$I(\chi_1,\chi_2)$ 典型地已经是不可约表示；而当两者相等时，它恰裂成一维表示与 Steinberg 扭曲的直和。完整的不可约性判别需要更系统的诱导表示理论，本章不再展开；这里只记录主系列和 Steinberg 作为 $\GL_2(\F_q)$ 表示论中最基本的两类对象。
\end{remark}

\section{Hecke 算子作为表示论对象}

上一节中的关键词是：某个代数（例如群代数，或由双陪集生成的代数）作用在一个向量空间上，从而把该向量空间变成一个模。Hecke 算子正是这一思想在模形式空间上的具体体现。

\subsection*{抽象 Hecke 代数}

\begin{definition}
设 $G$ 是一个群，$K\subseteq G$ 是一个子群。若双陪集 $KgK$ 的乘法在适当有限性条件下是良定义的，则以这些双陪集为基的复向量空间可赋予卷积乘法，得到\emph{Hecke 代数}，记作 $\mathcal{H}(G,K)$。

在有限群情形中，更具体地，$\mathcal{H}(G,K)$ 可定义为所有双 $K$-不变函数
\[
\varphi:G\to \C
\]
所成的向量空间，乘法为卷积
\[
(\varphi_1*\varphi_2)(g)=\sum_{x\in G}\varphi_1(x)\varphi_2(x^{-1}g).
\]
\end{definition}

这种代数通常不是交换的；但在很多与模形式相关的重要情形中，它恰恰是交换的。交换性意味着表示空间可以尝试分解为共同特征空间，这正是 Hecke 特征模形式出现的根源。

\subsection*{Hecke 算子在模形式空间上的作用}

设 $\Gamma\subseteq \SL_2(\Z)$ 为一个同余子群，$k\ge 0$。若
\[
\alpha=\mat{a&b\\c&d}\in \GL_2^+(\Q),
\]
定义权 $k$ 的 slash 作用
\[
(f|_k\alpha)(z)=\det(\alpha)^{k/2}(cz+d)^{-k}f\left(\frac{az+b}{cz+d}\right).
\]
当 $\Gamma\alpha\Gamma$ 是有限个右陪集之并
\[
\Gamma\alpha\Gamma=\bigsqcup_{i=1}^r \Gamma\alpha_i
\]
时，相应的双陪集算子定义为
\[
T_{\Gamma\alpha\Gamma}f=\sum_{i=1}^r f|_k\alpha_i.
\]
它只依赖于双陪集本身，而不依赖于右陪集代表元的选择。于是模形式空间 $\Mk_k(\Gamma)$ 成为某个 Hecke 代数的表示空间。

在最经典的情形 $\Gamma=\SL_2(\Z)$ 中，记标准 Hecke 算子为 $T_n$。若
\[
f(z)=\sum_{m=0}^{\infty}a_mq^m,\qquad q=e^{2\pi i z},
\]
则 $T_n f$ 的 $q$-展开公式为
\[
T_n f=\sum_{m=0}^{\infty}\left(\sum_{d\mid (m,n)} d^{k-1}a_{mn/d^2}\right)q^m.
\]
这是最适合进行代数运算的形式。

\begin{proposition}
设 $f(z)=\sum_{m\ge 0}a_mq^m\in \Mk_k(\SL_2(\Z))$。则
\begin{enumerate}[label=\textnormal{(\arabic*)}]
  \item 若 $(m,n)=1$，则 $T_mT_n=T_{mn}$；
  \item 对任意素数 $p$ 与整数 $r\ge 1$，有
  \[
  T_pT_{p^r}=T_{p^{r+1}}+p^{k-1}T_{p^{r-1}}.
  \]
\end{enumerate}
特别地，所有 Hecke 算子两两可交换。
\end{proposition}

\begin{proof}
记 $T_nf=\sum_{s\ge 0}b_sq^s$，其中
\[
b_s=\sum_{d\mid (s,n)}d^{k-1}a_{sn/d^2}.
\]

先证 (1)。设 $(m,n)=1$，则 $T_mT_nf$ 的 $q^r$ 系数为
\[
\sum_{d\mid (r,m)}d^{k-1}b_{rm/d^2}
=\sum_{d\mid (r,m)}\sum_{e\mid (rm/d^2,n)}d^{k-1}e^{k-1}a_{rmn/(d^2e^2)}.
\]
因为 $(m,n)=1$，且 $e\mid n$、$d\mid m$，故 $(d,e)=1$。令 $u=de$，则 $u\mid (r,mn)$，且每个这样的 $u$ 都唯一分解成 $d\mid m$ 与 $e\mid n$ 的乘积。于是上式化为
\[
\sum_{u\mid (r,mn)}u^{k-1}a_{rmn/u^2},
\]
这正是 $T_{mn}f$ 的 $q^r$ 系数。故 $T_mT_n=T_{mn}$。

再证 (2)。把 $T_{p^r}f$ 的系数记为
\[
b_s=\sum_{j=0}^{\min(v_p(s),r)}p^{j(k-1)}a_{sp^r/p^{2j}}.
\]
由于 $T_pf$ 的公式为
\[
T_pf=\sum_{m\ge 0}\bigl(a_{mp}+p^{k-1}\mathbf{1}_{p\mid m}a_{m/p}\bigr)q^m,
\]
故 $T_pT_{p^r}f$ 的 $q^m$ 系数为
\[
c_m=b_{mp}+p^{k-1}\mathbf{1}_{p\mid m}b_{m/p}.
\]
将 $b_{mp}$ 与 $b_{m/p}$ 的定义代入并整理求和指标，可得
\[
c_m=
\sum_{i=0}^{\min(v_p(m),r+1)}p^{i(k-1)}a_{mp^{r+1}/p^{2i}}
+ p^{k-1}\sum_{i=0}^{\min(v_p(m),r-1)}p^{i(k-1)}a_{mp^{r-1}/p^{2i}}.
\]
第一项恰是 $T_{p^{r+1}}f$ 的 $q^m$ 系数，第二项恰是 $p^{k-1}T_{p^{r-1}}f$ 的 $q^m$ 系数，因此得到所需关系。

现在对任意正整数按素因子分解，并利用 (1)、(2) 递归地把任意 $T_n$ 写成素数幂算子的多项式。因为这些关系对调换乘法顺序没有影响，所以 $T_mT_n=T_nT_m$，即 Hecke 算子两两可交换。
\end{proof}

\begin{definition}
设 $\mathbb{T}_k$ 为由全部 Hecke 算子生成的子代数
\[
\mathbb{T}_k\subseteq \End_{\C}(\Mk_k(\Gamma)).
\]
它称为权 $k$ 的\emph{Hecke 代数}。若非零模形式 $f\in \Mk_k(\Gamma)$ 满足对每个 $n\ge 1$ 都有
\[
T_nf=\lambda_n f,
\]
则称 $f$ 为一个\emph{Hecke 特征向量}或\emph{特征模形式}（eigenform）。
\end{definition}

交换代数作用在有限维向量空间上时，共同特征向量总是存在。

\begin{proposition}
设 $A$ 是代数闭域 $F$ 上的交换代数，$V$ 是一个非零有限维 $A$-模。则 $V$ 中存在非零向量 $v$，使得每个 $a\in A$ 都满足
\[
a\cdot v=\lambda(a)v
\]
对某个标量 $\lambda(a)\in F$ 成立。换言之，$V$ 至少有一个共同特征向量。
\end{proposition}

\begin{proof}
对 $\dim V$ 作归纳。若 $\dim V=1$，显然成立。设命题对较小维数成立。

若 $A$ 中每个元素在 $V$ 上都按标量作用，则任取 $0\neq v\in V$ 即为共同特征向量。否则可取某个 $a\in A$，使其在线性空间 $V$ 上不是标量变换。由于 $F$ 代数闭，$a:V\to V$ 至少有一个特征值 $\lambda$，相应特征子空间
\[
E=\Ker(a-\lambda\Id)
\]
非零。因为 $a$ 不是标量，所以 $E\neq V$，从而 $\dim E<\dim V$。

对任意 $b\in A$ 与 $v\in E$，因为 $ab=ba$，有
\[
(a-\lambda\Id)(b\cdot v)=b\cdot (a-\lambda\Id)(v)=0,
\]
故 $b\cdot v\in E$。因此 $E$ 是一个非零真 $A$-子模。由归纳假设，$E$ 中存在共同特征向量，它当然也是 $V$ 中的共同特征向量。命题得证。
\end{proof}

\begin{proposition}
设
\[
f(q)=\sum_{n=0}^{\infty}a_nq^n\in \Mk_k(\SL_2(\Z))
\]
是 Hecke 特征向量，满足 $a_1\neq 0$，且
\[
T_nf=\lambda_n f\qquad (n\ge 1).
\]
则
\[
\lambda_n=\frac{a_n}{a_1}.
\]
特别地，若 $f$ 规范化为 $a_1=1$，则每个 Hecke 特征值恰等于相应 Fourier 系数：$\lambda_n=a_n$。
\end{proposition}

\begin{proof}
由 $T_nf=\lambda_nf$，比较两边 $q$ 展开的 $q^1$ 系数即可。根据 Hecke 算子的公式，$T_nf$ 的 $q^1$ 系数为
\[
\sum_{d\mid (1,n)}d^{k-1}a_{n/d^2}=a_n.
\]
而 $\lambda_nf$ 的 $q^1$ 系数为 $\lambda_n a_1$。因此
\[
a_n=\lambda_n a_1,
\]
即
\[
\lambda_n=\frac{a_n}{a_1}.
\]
\end{proof}

\begin{remark}[与自守表示的联系：预告]
现代观点认为，模形式不只是某个函数空间中的向量，更应看作 $\GL_2$ 的自守表示中的向量。Hecke 代数在模形式空间上的作用，对应于局部双陪集代数在自守表示上的作用；未分歧处的局部 Hecke 代数是交换代数，因此在不可约自守表示上按 Schur 引理只能以标量方式作用。这些标量正是我们熟悉的 Hecke 特征值。

更具体地说，一个规范化 Hecke 本征 cusp form $f$ 通常对应一个不可约自守表示 $\pi_f=\bigotimes_v \pi_{f,v}$。对于几乎所有素数 $p$，局部分量 $\pi_{f,p}$ 是未分歧的，而 $T_p$ 的特征值 $a_p(f)$ 记录了该局部分量的 Satake 参数。换句话说，本章中“交换代数作用于有限维空间”的初等图景，正是后续自守表示论的有限维影子。
\end{remark}

\section*{习题}
\addcontentsline{toc}{section}{习题}

下面共有 $60$ 题，按难度分为基础（\basic）、进阶（\medium）与挑战（\hard）三级。

\subsection*{基础题（共 25 题）}

\begin{enumerate}[label={\textbf{习题8.\arabic*.}}, wide=0pt, leftmargin=3.2em]
\item \basic 设 $G$ 为任意群，验证映射 $\rho(g)=\Id_F$ 的确给出一个表示。

\item \basic 求循环群 $C_2$ 在 $\C$ 上的全部不可约表示，并写出它们的特征标。

\item \basic 写出 $S_3$ 在 $\C^3$ 上置换坐标作用时，换位 $(12)$ 与三循环 $(123)$ 的矩阵。

\item \basic 验证 $\sgn:S_n\to \{\pm 1\}$ 是群同态。

\item \basic 写出 $C_3=\langle r\rangle$ 的左正则表示中 $\rho(r)$ 在基 $e_1,e_r,e_{r^2}$ 下的矩阵。

\item \basic 设 $\rho:G\to \GL(V)$ 为表示，$T\in \GL(V)$。验证 $\rho'(g)=T\rho(g)T^{-1}$ 也是表示，并与 $\rho$ 等价。

\item \basic 在 $S_3$ 的三维置换表示中，证明子空间
\[
\{(x,x,x):x\in \C\}
\]
与
\[
\{(x_1,x_2,x_3):x_1+x_2+x_3=0\}
\]
都是子表示。

\item \basic 设 $V=W_1\oplus W_2$ 是表示的直和，证明 $\chi_V=\chi_{W_1}+\chi_{W_2}$。

\item \basic 对群 $C_3$ 在 $\C^3$ 上的正则表示，选取子表示 $W=\C(e_1+e_r+e_{r^2})$，按照 Maschke 定理的平均构造写出到 $W$ 的 $G$-等变投影。

\item \basic 设 $G=C_p$，系数域为 $\F_p$，表示由
\[
\rho(g)=\mat{1&1\\0&1}
\]
给出。验证 $F_p\mat{1\\0}$ 是子表示，但没有不变补空间。

\item \basic 证明任意一维表示都是不可约的。

\item \basic 计算正则表示的特征标 $\chi_{\mathrm{reg}}$。

\item \basic 设 $G$ 作用在有限集合 $X$ 上，证明置换表示 $\C[X]$ 的特征标值 $\chi(g)$ 等于 $g$ 在 $X$ 上的不动点个数。

\item \basic 写出循环群 $C_4$ 的特征标表。

\item \basic 写出循环群 $C_5$ 的特征标表。

\item \basic 根据正文中的信息补全 $S_3$ 的特征标表，并直接检查第一正交关系。

\item \basic 写出 $D_4$ 的共轭类，并说明为什么它恰有 $5$ 个不可约复表示。

\item \basic 设 $g=\mat{a&0\\0&d}\in \GL_2(F)$。计算 $g$ 在 $\mathrm{Sym}^2(F^2)$ 上的作用矩阵及其迹。

\item \basic 证明 $\dim \mathrm{Sym}^k(F^2)=k+1$。

\item \basic 利用 Hecke 算子的 $q$-展开公式证明：若 $(m,n)=1$，则 $T_mT_n=T_{mn}$。

\item \basic 设 $A,B\in M_2(\C)$ 可交换。证明它们至少有一个公共特征向量。

\item \basic 设 $C_4$ 作用在正方形顶点集上，写出对应置换表示的特征标。

\item \basic 求 $S_3$ 的类函数空间维数。

\item \basic 验证循环群 $C_n$ 的特征标 $\chi_j(t^m)=\zeta_n^{jm}$ 两两正交。

\item \basic 计算 $S_3$ 的表示 $\sgn\otimes \mathrm{Std}$ 的特征标，并判断它与哪一个不可约表示等价。
\end{enumerate}

\subsection*{进阶题（共 25 题）}

\begin{enumerate}[label={\textbf{习题8.\arabic*.}}, wide=0pt, leftmargin=3.2em]
\setcounter{enumi}{25}

\item \medium 证明：子表示 $W\subseteq V$ 有不变补空间，当且仅当存在像为 $W$ 的 $G$-等变投影 $P:V\to V$。

\item \medium 设 $V_1,\dots,V_r$ 为 $G$ 的全部不可约复表示。利用正则表示证明
\[
|G|=\sum_{i=1}^r (\dim V_i)^2.
\]

\item \medium 证明：有限阿贝尔群的任意不可约复表示都是一维的。

\item \medium 设 $V$ 是不可约复表示，证明群代数中心 $Z(\C[G])$ 在 $V$ 上按标量作用。

\item \medium 设 $V,W$ 为表示，证明 $V\otimes W$ 上的自然作用确实定义了一个表示，并验证其特征标满足
\[
\chi_{V\otimes W}=\chi_V\chi_W.
\]

\item \medium 设 $S_3$ 作用在有序对集合
\[
X=\{(i,j):1\le i,j\le 3,\ i\neq j\}
\]
上。求置换表示 $\C[X]$ 的特征标，并将其分解为不可约表示直和。

\item \medium 将 $S_3$ 的正则表示分解为不可约表示直和。

\item \medium 仅利用共轭类与平方和公式，重新推导 $D_4$ 的不可约表示维数只能是 $1,1,1,1,2$。

\item \medium 用第二正交关系验证 $S_3$ 特征标表的列正交性。

\item \medium 对一般循环群 $C_n$，证明其特征标表就是离散 Fourier 矩阵。

\item \medium 设 $H\triangleleft G$ 且 $[G:H]=2$。证明由商映射 $G\to G/H\cong C_2$ 得到一个非平凡一维表示，并说明对应的置换表示如何分解为两个一维表示。

\item \medium 设 $\mathrm{Std}$ 为 $S_3$ 的二维标准表示。利用特征标分解 $\mathrm{Sym}^2(\mathrm{Std})$。

\item \medium 证明：若一个有限群在有限集合 $X$ 上 $2$-重传递，则置换表示 $\C[X]$ 的自同态环维数为 $2$。

\item \medium 设 $g=\mathrm{diag}(a,d)\in \GL_2(F)$。求 $g$ 在 $\mathrm{Sym}^k(F^2)$ 上的全部特征值与特征标。

\item \medium 证明 $\mathrm{St}$ 的维数为 $q$，并验证
\[
\C[\mathbf{P}^1(\F_q)]\cong \mathbf{1}\oplus \mathrm{St}.
\]

\item \medium 证明
\[
I(\chi,\chi)\cong (\chi\circ\det)\otimes I(1,1).
\]

\item \medium 设 $A$ 是代数闭域上的交换代数，$V$ 是有限维 $A$-模。用归纳法证明 $V$ 中存在共同特征向量。

\item \medium 设 $f(q)=\sum_{n\ge 1}a_nq^n$ 是规范化 Hecke 本征模形式。证明其 Hecke 特征值满足 $\lambda_n=a_n$。

\item \medium 利用 Hecke 关系证明：若 $(m,n)=1$，则规范化本征模形式的 Fourier 系数满足 $a_{mn}=a_ma_n$。

\item \medium 求 $\sgn\otimes \sgn$、$\sgn\otimes \mathrm{Std}$、$\mathrm{Std}\otimes \mathrm{Std}$ 的 $S_3$ 特征标，并将它们分解为不可约表示直和。

\item \medium 证明：若两个复表示特征标相同，则它们等价。

\item \medium 求 $D_4$ 在四个顶点上的置换表示的特征标，并将其分解为不可约表示直和。

\item \medium 设 $V,W$ 为复表示，证明
\[
\dim \Hom_G(V,W)=\inner{\chi_V}{\chi_W}.
\]

\item \medium 证明同构
\[
\Hom_G(V,W)\cong (V^{\vee}\otimes W)^G,
\]
并解释它与 Schur 引理的关系。

\item \medium 证明：若 $V$ 是不可约复表示，则其对偶表示 $V^{\vee}$ 的特征标满足
\[
\chi_{V^{\vee}}(g)=\overline{\chi_V(g)}.
\]
\end{enumerate}

\subsection*{挑战题（共 10 题）}

\begin{enumerate}[label={\textbf{习题8.\arabic*.}}, wide=0pt, leftmargin=3.2em]
\setcounter{enumi}{50}

\item \hard 设 $H\subseteq G$，$W$ 是 $H$ 的表示。按函数模型定义诱导表示 $\mathrm{Ind}_H^G W$，并证明
\[
\dim \mathrm{Ind}_H^G W=[G:H]\dim W.
\]

\item \hard Frobenius 互反律断言
\[
\Hom_G(\mathrm{Ind}_H^G W,V)\cong \Hom_H(W,\mathrm{Res}_H^G V).
\]
请先准确陈述该定理，然后在 $H\triangleleft G$ 且 $W=\mathbf{1}_H$ 的情形下验证它。

\item \hard 证明：有限阿贝尔群 $A$ 的不可约特征标构成的群 $\widehat{A}$ 与 $A$ 同构（非规范同构即可）。

\item \hard 设 $p$ 为奇素数。对 $G=\GL_2(\F_p)$，证明 $G/B\cong \mathbf{P}^1(\F_p)$，并由此重新得到主系列表示 $I(\chi_1,\chi_2)$ 的维数是 $p+1$。

\item \hard 设 $G=\GL_2(\F_p)$ 在 $\mathbf{P}^1(\F_p)$ 上作用，记其置换特征标为 $\pi$。求 $\pi(g)$ 与 $g$ 在射影直线上的不动点个数之间的关系，并据此讨论恒等元、非平凡幂零元、对角化元三类情形下的取值。

\item \hard 证明
\[
\mathrm{Ind}_B^{\GL_2(\F_q)}\mathbf{1}\cong \mathbf{1}\oplus \mathrm{St}.
\]
并说明这正是 Steinberg 表示最基本的构造方式。

\item \hard 设有限群 $G$ 作用在集合 $X$ 上，$x\in X$ 的稳定子为 $H$。利用 Frobenius 互反律证明：平凡表示在置换表示 $\C[X]$ 中出现的重数等于轨道个数；在传递作用情形下它恰出现一次。

\item \hard 设 $f=\sum_{n\ge 1}a_nq^n$ 是权 $k$ 的规范化 Hecke 本征 cusp form。利用
\[
T_pT_{p^r}=T_{p^{r+1}}+p^{k-1}T_{p^{r-1}}
\]
证明递推关系
\[
a_{p^{r+1}}=a_pa_{p^r}-p^{k-1}a_{p^{r-1}}.
\]

\item \hard 在上一题基础上证明局部 Euler 因子恒等式
\[
\sum_{r\ge 0}a_{p^r}X^r=\frac{1}{1-a_pX+p^{k-1}X^2}.
\]
并说明这为何暗示 Hecke 特征值应由某个二维表示控制。

\item \hard 阅读本节关于自守表示的预告，说明下列口号的精确含义：
“模形式的 Hecke 本征值就是某个 $\GL_2$ 表示在未分歧局部 Hecke 代数上的特征标量。”
要求给出你自己的分层解释：先说有限维线性代数图景，再说群表示图景，最后说自守表示图景。
\end{enumerate}

\section*{习题解答}
\begin{enumerate}[label=\textbf{\arabic*.}, leftmargin=*, itemsep=1em]
\item \textbf{解答.} 记 $\rho(g)=\Id_F$。则
\[
\rho(1)=\Id_F,\qquad \rho(gh)=\Id_F=\Id_F\Id_F=\rho(g)\rho(h)
\]
对任意 $g,h\in G$ 都成立，故 $\rho:G\to \GL(F)$ 是群同态，因此确为一个表示。

\item \textbf{解答.} 设 $C_2=\langle t\mid t^2=1\rangle$。因 $C_2$ 阿贝尔，其不可约复表示全为一维。故只需令
\[
\rho(t)=\lambda\in \C^\times,\qquad \lambda^2=1,
\]
于是 $\lambda=\pm1$。故全部不可约表示只有两个：
\[
\mathbf{1}:t\mapsto 1,\qquad \varepsilon:t\mapsto -1.
\]
其特征标分别为
\[
\chi_{\mathbf{1}}(1)=1,\ \chi_{\mathbf{1}}(t)=1;\qquad
\chi_\varepsilon(1)=1,\ \chi_\varepsilon(t)=-1.
\]

\item \textbf{解答.} 在标准基 $e_1,e_2,e_3$ 下，置换表示满足 $\sigma\cdot e_i=e_{\sigma(i)}$。故
\[
(12):\ e_1\leftrightarrow e_2,\ e_3\mapsto e_3,
\]
所以
\[
\rho((12))=\mat{0&1&0\\1&0&0\\0&0&1}.
\]
又
\[
(123):\ e_1\mapsto e_2,\ e_2\mapsto e_3,\ e_3\mapsto e_1,
\]
所以
\[
\rho((123))=\mat{0&0&1\\1&0&0\\0&1&0}.
\]

\item \textbf{解答.} 置换的符号可定义为任一换位分解长度的奇偶性：
\[
\sgn(\sigma)=(-1)^{\ell(\sigma)}.
\]
若 $\sigma=\tau_1\cdots \tau_r,\ \pi=\eta_1\cdots \eta_s$ 为换位分解，则
\[
\sigma\pi=\tau_1\cdots \tau_r\eta_1\cdots \eta_s
\]
是一个换位分解，其长度奇偶性为 $r+s$，故
\[
\sgn(\sigma\pi)=(-1)^{r+s}=(-1)^r(-1)^s=\sgn(\sigma)\sgn(\pi).
\]
因此 $\sgn:S_n\to\{\pm1\}$ 为群同态。

\item \textbf{解答.} 左正则表示满足
\[
\rho(r)e_1=e_r,\qquad \rho(r)e_r=e_{r^2},\qquad \rho(r)e_{r^2}=e_1.
\]
故在基 $e_1,e_r,e_{r^2}$ 下，
\[
\rho(r)=\mat{0&0&1\\1&0&0\\0&1&0}.
\]

\item \textbf{解答.} 设
\[
\rho'(g)=T\rho(g)T^{-1}.
\]
则
\[
\rho'(gh)=T\rho(gh)T^{-1}=T\rho(g)\rho(h)T^{-1}
=(T\rho(g)T^{-1})(T\rho(h)T^{-1})
=\rho'(g)\rho'(h),
\]
且 $\rho'(1)=T\Id T^{-1}=\Id$，故 $\rho'$ 是表示。又
\[
T\rho(g)=\rho'(g)T
\]
对一切 $g$ 成立，所以 $T$ 是交织同构，故 $\rho'\cong \rho$。

\item \textbf{解答.} 设
\[
U=\{(x,x,x):x\in\C\},\qquad
W=\{(x_1,x_2,x_3):x_1+x_2+x_3=0\}.
\]
对任意 $\sigma\in S_3$，若 $v=(x,x,x)\in U$，则 $\sigma\cdot v$ 只是重排坐标，仍为 $(x,x,x)$，故 $U$ 不变。若 $w=(x_1,x_2,x_3)\in W$，则
\[
(\sigma\cdot w)_1+(\sigma\cdot w)_2+(\sigma\cdot w)_3=x_1+x_2+x_3=0,
\]
故 $\sigma\cdot w\in W$。因此 $U,W$ 都是子表示。

\item \textbf{解答.} 对任意 $g\in G$，在分解 $V=W_1\oplus W_2$ 下，
\[
\rho_V(g)=\mat{\rho_{W_1}(g)&0\\0&\rho_{W_2}(g)}.
\]
故迹满足
\[
\chi_V(g)=\Tr(\rho_V(g))
=\Tr(\rho_{W_1}(g))+\Tr(\rho_{W_2}(g))
=\chi_{W_1}(g)+\chi_{W_2}(g).
\]
于是 $\chi_V=\chi_{W_1}+\chi_{W_2}$。

\item \textbf{解答.} 记
\[
u=e_1+e_r+e_{r^2},\qquad W=\C u.
\]
先取线性投影 $P_0$ 为
\[
P_0(e_1)=u,\qquad P_0(e_r)=P_0(e_{r^2})=0.
\]
按 Maschke 平均构造
\[
P=\frac13\sum_{g\in C_3}\rho(g)P_0\rho(g)^{-1}.
\]
对
\[
v=x_1e_1+x_re_r+x_{r^2}e_{r^2}
\]
直接计算得
\[
P(v)=\frac{x_1+x_r+x_{r^2}}{3}\,(e_1+e_r+e_{r^2}).
\]
这显然像在 $W$ 中，且 $P(u)=u$，故是到 $W$ 的 $G$-等变投影。

\item \textbf{解答.} 设
\[
\rho(g)=\mat{1&1\\0&1},\qquad W=\F_p\mat{1\\0}.
\]
则
\[
\rho(g)\mat{1\\0}=\mat{1\\0}\in W,
\]
故 $W$ 是子表示。若存在不变补空间 $U$，则因 $\dim V=2$，$U$ 必为一维，且可写成
\[
U=\F_p\mat{a\\1}
\]
（否则 $U=W$）。若 $U$ 不变，则
\[
\rho(g)\mat{a\\1}=\mat{a+1\\1}\in U,
\]
故存在 $\lambda\in\F_p$ 使
\[
\mat{a+1\\1}=\lambda\mat{a\\1}.
\]
比较第二坐标得 $\lambda=1$，再比较第一坐标得 $a+1=a$，矛盾。故 $W$ 没有不变补空间。

\item \textbf{解答.} 设 $\dim V=1$。则 $V$ 的子空间只有 $0$ 与 $V$。它们显然都是不变子空间，因此任意一维表示都不可约。

\item \textbf{解答.} 设 $\rho_{\mathrm{reg}}$ 为左正则表示，基为 $\{e_h:h\in G\}$。对 $g\in G$，
\[
\rho_{\mathrm{reg}}(g)e_h=e_{gh}.
\]
若 $g=1$，则 $\rho_{\mathrm{reg}}(1)=\Id$，故
\[
\chi_{\mathrm{reg}}(1)=|G|.
\]
若 $g\neq1$，则方程 $gh=h$ 无解，所以基向量全被置换而无不动点，矩阵对角元全为 $0$，故
\[
\chi_{\mathrm{reg}}(g)=0.
\]
因此
\[
\chi_{\mathrm{reg}}(g)=
\begin{cases}
|G|,& g=1,\\
0,& g\neq1.
\end{cases}
\]

\item \textbf{解答.} 置换表示 $\C[X]$ 以 $\{e_x:x\in X\}$ 为基，且
\[
\rho(g)e_x=e_{g\cdot x}.
\]
故 $\rho(g)$ 的矩阵是一个置换矩阵。其对角线上第 $x$ 个位置为 $1$ 当且仅当 $g\cdot x=x$。因此
\[
\chi(g)=\Tr(\rho(g))=\#\{x\in X:g\cdot x=x\},
\]
即等于 $g$ 在 $X$ 上的不动点个数。

\item \textbf{解答.} 设 $C_4=\langle r\rangle$，取 $i=\sqrt{-1}$。其四个不可约特征标为
\[
\chi_j(r^m)=i^{jm}\qquad (j=0,1,2,3).
\]
特征标表为
\[
\begin{array}{c|cccc}
 & 1 & r & r^2 & r^3\\
\hline
\chi_0 & 1 & 1 & 1 & 1\\
\chi_1 & 1 & i & -1 & -i\\
\chi_2 & 1 & -1 & 1 & -1\\
\chi_3 & 1 & -i & -1 & i
\end{array}
\]

\item \textbf{解答.} 设 $C_5=\langle r\rangle$，取本原五次单位根 $\zeta=\zeta_5$。其五个不可约特征标为
\[
\chi_j(r^m)=\zeta^{jm}\qquad (j=0,1,2,3,4).
\]
特征标表为
\[
\begin{array}{c|ccccc}
 & 1 & r & r^2 & r^3 & r^4\\
\hline
\chi_0 & 1 & 1 & 1 & 1 & 1\\
\chi_1 & 1 & \zeta & \zeta^2 & \zeta^3 & \zeta^4\\
\chi_2 & 1 & \zeta^2 & \zeta^4 & \zeta & \zeta^3\\
\chi_3 & 1 & \zeta^3 & \zeta & \zeta^4 & \zeta^2\\
\chi_4 & 1 & \zeta^4 & \zeta^3 & \zeta^2 & \zeta
\end{array}
\]

\item \textbf{解答.} $S_3$ 的共轭类与大小分别为
\[
\{1\}\ (1),\qquad \{(12),(13),(23)\}\ (3),\qquad \{(123),(132)\}\ (2).
\]
特征标表为
\[
\begin{array}{c|ccc}
 & 1 & (12) & (123)\\
\hline
\mathbf{1} & 1 & 1 & 1\\
\sgn & 1 & -1 & 1\\
\mathrm{Std} & 2 & 0 & -1
\end{array}
\]
直接检验第一正交关系：
\[
\inner{\mathbf{1}}{\sgn}
=\frac16(1\cdot1+3\cdot1\cdot(-1)+2\cdot1\cdot1)=0,
\]
\[
\inner{\mathbf{1}}{\mathrm{Std}}
=\frac16(1\cdot2+3\cdot1\cdot0+2\cdot1\cdot(-1))=0,
\]
\[
\inner{\sgn}{\mathrm{Std}}
=\frac16(1\cdot2+3\cdot(-1)\cdot0+2\cdot1\cdot(-1))=0.
\]
且
\[
\inner{\mathbf{1}}{\mathbf{1}}=\inner{\sgn}{\sgn}=1,\qquad
\inner{\mathrm{Std}}{\mathrm{Std}}=\frac16(4+0+2)=1.
\]
故各不可约特征标两两正交且范数为 $1$。

\item \textbf{解答.} 由正文，
\[
D_4=\langle r,s\mid r^4=s^2=1,\ srs=r^{-1}\rangle
\]
的共轭类为
\[
\{1\},\quad \{r^2\},\quad \{r,r^3\},\quad \{s,sr^2\},\quad \{sr,sr^3\}.
\]
有限群的不可约复表示个数等于其共轭类个数，故 $D_4$ 恰有 $5$ 个不可约复表示。

\item \textbf{解答.} 在基 $X^2,XY,Y^2$ 下，$g=\diag(a,d)$ 的作用为
\[
g\cdot X=aX,\qquad g\cdot Y=dY.
\]
故
\[
g\cdot X^2=a^2X^2,\qquad g\cdot (XY)=ad\,XY,\qquad g\cdot Y^2=d^2Y^2.
\]
因此作用矩阵为
\[
\mat{a^2&0&0\\0&ad&0\\0&0&d^2},
\]
其迹为
\[
a^2+ad+d^2.
\]

\item \textbf{解答.} 将 $\mathrm{Sym}^k(F^2)$ 视为二元 $k$ 次齐次多项式空间。单项式
\[
X^k,\ X^{k-1}Y,\ \dots,\ XY^{k-1},\ Y^k
\]
共有 $k+1$ 个，且线性无关并张成全体齐次多项式，故它们构成一组基。因此
\[
\dim \mathrm{Sym}^k(F^2)=k+1.
\]

\item \textbf{解答.} 设
\[
f(q)=\sum_{r\ge0}a_rq^r,\qquad T_nf=\sum_{s\ge0}b_sq^s,
\]
则
\[
b_s=\sum_{d\mid(s,n)}d^{k-1}a_{sn/d^2}.
\]
若 $(m,n)=1$，则 $T_mT_nf$ 的 $q^r$ 系数为
\[
\sum_{d\mid(r,m)}d^{k-1}b_{rm/d^2}
=\sum_{d\mid(r,m)}\sum_{e\mid(rm/d^2,n)}d^{k-1}e^{k-1}a_{rmn/(d^2e^2)}.
\]
因 $d\mid m,\ e\mid n,\ (m,n)=1$，故 $(d,e)=1$。令 $u=de$，则 $u\mid(r,mn)$，且每个这样的 $u$ 唯一写成 $de$。于是上式化为
\[
\sum_{u\mid(r,mn)}u^{k-1}a_{rmn/u^2},
\]
这正是 $T_{mn}f$ 的 $q^r$ 系数。故
\[
T_mT_n=T_{mn}.
\]

\item \textbf{解答.} 先取 $A$ 的一个特征值 $\lambda$，设其特征子空间
\[
E_\lambda=\Ker(A-\lambda I)\neq0.
\]
因 $AB=BA$，对任意 $v\in E_\lambda$，
\[
A(Bv)=B(Av)=B(\lambda v)=\lambda Bv,
\]
故 $Bv\in E_\lambda$，即 $E_\lambda$ 在 $B$ 下不变。若 $\dim E_\lambda=1$，则其任一非零向量自动是 $B$ 的特征向量；若 $\dim E_\lambda=2$，则 $E_\lambda=\C^2$，而 $B$ 在复数域上必有特征向量。故总存在 $v\neq0$，使
\[
Av=\lambda v,\qquad Bv=\mu v.
\]
即 $A,B$ 至少有一个公共特征向量。

\item \textbf{解答.} $C_4$ 在正方形四顶点上的旋转作用中，恒等元固定全部 $4$ 个顶点，而 $r,r^2,r^3$ 都不固定任何顶点。故对应置换表示的特征标为
\[
\chi(1)=4,\qquad \chi(r)=\chi(r^2)=\chi(r^3)=0.
\]

\item \textbf{解答.} 类函数空间的维数等于共轭类个数。$S_3$ 有三个共轭类：
\[
\{1\},\qquad \{(12),(13),(23)\},\qquad \{(123),(132)\}.
\]
故 $S_3$ 的类函数空间维数为
\[
3.
\]

\item \textbf{解答.} 取类函数内积。对 $0\le i,j\le n-1$，
\[
\inner{\chi_i}{\chi_j}
=\frac1n\sum_{m=0}^{n-1}\chi_i(t^m)\overline{\chi_j(t^m)}
=\frac1n\sum_{m=0}^{n-1}\zeta_n^{(i-j)m}.
\]
若 $i=j$，则和为 $n$，故内积为 $1$。若 $i\neq j$，则 $\zeta_n^{i-j}\neq1$，故几何级数求和得
\[
\sum_{m=0}^{n-1}\zeta_n^{(i-j)m}
=\frac{1-\zeta_n^{(i-j)n}}{1-\zeta_n^{i-j}}=0.
\]
故
\[
\inner{\chi_i}{\chi_j}=\delta_{ij},
\]
即它们两两正交。

\item \textbf{解答.} 特征标逐点相乘：
\[
\chi_{\sgn\otimes\mathrm{Std}}
=\chi_{\sgn}\chi_{\mathrm{Std}}
=(1,-1,1)\cdot(2,0,-1)=(2,0,-1).
\]
这正是 $\mathrm{Std}$ 的特征标，因此
\[
\sgn\otimes\mathrm{Std}\cong \mathrm{Std}.
\]

\item \textbf{解答.} 若 $V=W\oplus U$，其中 $W,U$ 都是子表示，则沿 $U$ 到 $W$ 的投影 $P$ 满足
\[
P^2=P,\qquad \Ima(P)=W.
\]
对 $v=w+u$（$w\in W,u\in U$）有
\[
P(gv)=P(gw+gu)=gw=gP(v),
\]
故 $P$ 为 $G$-等变投影。

反之，若存在 $G$-等变投影 $P:V\to V$，且 $\Ima(P)=W$，令
\[
U=\Ker(P).
\]
由 $P^2=P$ 得
\[
V=\Ima(P)\oplus\Ker(P)=W\oplus U.
\]
又若 $u\in U$，则
\[
P(gu)=gP(u)=0,
\]
故 $gu\in U$，所以 $U$ 也是子表示。于是 $W$ 有不变补空间。

\item \textbf{解答.} 设 $V_1,\dots,V_r$ 为全部不可约复表示，特征标分别为 $\chi_i$。正则表示分解为
\[
\C[G]\cong \bigoplus_{i=1}^r m_iV_i.
\]
其重数
\[
m_i=\inner{\chi_{\mathrm{reg}}}{\chi_i}
=\frac1{|G|}\sum_{g\in G}\chi_{\mathrm{reg}}(g)\overline{\chi_i(g)}
=\frac1{|G|}\chi_{\mathrm{reg}}(1)\chi_i(1)
=\dim V_i.
\]
故
\[
\C[G]\cong \bigoplus_{i=1}^r (\dim V_i)\,V_i.
\]
两边取维数即得
\[
|G|=\dim \C[G]=\sum_{i=1}^r (\dim V_i)^2.
\]

\item \textbf{解答.} 设 $\rho:G\to \GL(V)$ 为不可约复表示，且 $G$ 阿贝尔。则对任意 $g,h\in G$，
\[
\rho(g)\rho(h)=\rho(h)\rho(g).
\]
于是每个 $\rho(g)$ 都与所有群作用可交换。由 Schur 引理，
\[
\rho(g)=\lambda_g\Id_V.
\]
因此任意一维子空间都不变。若 $\dim V>1$，则存在非平凡真子空间，与不可约性矛盾。故
\[
\dim V=1.
\]

\item \textbf{解答.} 任取
\[
z=\sum_{g\in G}a_gg\in Z(\C[G]).
\]
在表示 $V$ 上它作用为
\[
\rho(z)=\sum_{g\in G}a_g\rho(g).
\]
对任意 $h\in G$，
\[
\rho(h)\rho(z)=\rho(hz)=\rho(zh)=\rho(z)\rho(h).
\]
故 $\rho(z)\in \End_G(V)$。由于 $V$ 不可约且系数域为 $\C$，Schur 引理给出
\[
\rho(z)=\lambda(z)\Id_V.
\]
所以 $Z(\C[G])$ 在 $V$ 上按标量作用。

\item \textbf{解答.} 在 $V\otimes W$ 上定义
\[
g\cdot(v\otimes w)=(g\cdot v)\otimes(g\cdot w),
\]
并线性延拓。则
\[
1\cdot(v\otimes w)=v\otimes w,
\]
且
\[
g\cdot(h\cdot(v\otimes w))
=g\cdot((hv)\otimes(hw))
=(gh)v\otimes(gh)w
=(gh)\cdot(v\otimes w),
\]
故这确实定义了一个表示。

若 $\rho_V(g)$ 的特征值为 $\lambda_1,\dots,\lambda_m$，$\rho_W(g)$ 的特征值为 $\mu_1,\dots,\mu_n$，则 $\rho_{V\otimes W}(g)$ 的特征值为全部 $\lambda_i\mu_j$，故
\[
\chi_{V\otimes W}(g)=\sum_{i,j}\lambda_i\mu_j
=\left(\sum_i\lambda_i\right)\left(\sum_j\mu_j\right)
=\chi_V(g)\chi_W(g).
\]

\item \textbf{解答.} 集合
\[
X=\{(i,j):1\le i,j\le3,\ i\neq j\}
\]
有 $6$ 个元素。置换特征标等于不动点个数：

恒等元固定全部 $6$ 个点，故 $\chi(1)=6$。

换位如 $(12)$ 只固定数字 $3$，但 $(3,3)\notin X$，故无不动点，所以 $\chi((12))=0$。

三循环如 $(123)$ 无固定数字，故也无不动点，所以 $\chi((123))=0$。

于是
\[
\chi_{\C[X]}=(6,0,0).
\]
与 $S_3$ 不可约特征标作内积：
\[
\inner{\chi}{\mathbf1}=\frac16\cdot6=1,\qquad
\inner{\chi}{\sgn}=\frac16\cdot6=1,\qquad
\inner{\chi}{\mathrm{Std}}=\frac16(6\cdot2)=2.
\]
故
\[
\C[X]\cong \mathbf1\oplus \sgn\oplus \mathrm{Std}^{\oplus2}.
\]
事实上，这个作用是自由传递的，所以它正是正则表示。

\item \textbf{解答.} 由上一题或由一般公式，正则表示中不可约表示 $V_i$ 的重数为 $\dim V_i$。对 $S_3$ 而言，不可约表示为
\[
\mathbf1,\qquad \sgn,\qquad \mathrm{Std},
\]
维数分别为 $1,1,2$。故
\[
\C[S_3]\cong \mathbf1\oplus \sgn\oplus \mathrm{Std}^{\oplus2}.
\]

\item \textbf{解答.} $D_4$ 有 $5$ 个共轭类，故有 $5$ 个不可约表示。设其维数为正整数
\[
d_1,\dots,d_5.
\]
由平方和公式，
\[
d_1^2+\cdots+d_5^2=|D_4|=8.
\]
因为每个 $d_i\ge1$，故
\[
d_1^2+\cdots+d_5^2\ge 1+1+1+1+1=5.
\]
离 $8$ 还差 $3$。把一个 $1^2$ 改成 $2^2$ 恰好增加 $3$；若有某个 $d_i\ge3$，则平方至少为 $9$，不可能。故唯一可能是
\[
1,1,1,1,2.
\]

\item \textbf{解答.} 对 $S_3$，共轭类代表可取
\[
1,\quad (12),\quad (123),
\]
其中心化子大小分别为
\[
6,\quad 2,\quad 3.
\]
由特征标表
\[
\begin{array}{c|ccc}
 & 1 & (12) & (123)\\
\hline
\mathbf{1} & 1 & 1 & 1\\
\sgn & 1 & -1 & 1\\
\mathrm{Std} & 2 & 0 & -1
\end{array}
\]
直接算得
\[
1^2+1^2+2^2=6,\qquad
1^2+(-1)^2+0^2=2,\qquad
1^2+1^2+(-1)^2=3,
\]
对应列的范数正好是中心化子大小。又
\[
1\cdot1+1\cdot(-1)+2\cdot0=0,
\]
\[
1\cdot1+1\cdot1+2\cdot(-1)=0,
\]
\[
1\cdot1+(-1)\cdot1+0\cdot(-1)=0.
\]
故各列两两正交，这正是第二正交关系。

\item \textbf{解答.} 设 $C_n=\langle t\rangle$，取本原 $n$ 次单位根 $\zeta_n$。其特征标表第 $i$ 行第 $j$ 列为
\[
\chi_i(t^j)=\zeta_n^{ij}\qquad (0\le i,j\le n-1).
\]
因此特征标表就是矩阵
\[
F_n=(\zeta_n^{ij})_{0\le i,j\le n-1},
\]
这正是离散 Fourier 矩阵；归一化后 $\frac1{\sqrt n}F_n$ 为酉矩阵。

\item \textbf{解答.} 因 $H\triangleleft G$ 且 $[G:H]=2$，商群 $G/H\cong C_2$。设
\[
\varepsilon:G\to G/H\cong C_2\hookrightarrow \{\pm1\}\subset \C^\times.
\]
则 $\varepsilon$ 是一个非平凡一维表示，且
\[
\varepsilon(g)=
\begin{cases}
1,& g\in H,\\
-1,& g\notin H.
\end{cases}
\]

对 $G/H=\{H,gH\}$ 的置换表示 $\C[G/H]$，取基
\[
u=e_H+e_{gH},\qquad v=e_H-e_{gH}.
\]
则 $u$ 张成平凡子表示，而
\[
x\cdot v=\varepsilon(x)v
\]
对任意 $x\in G$ 成立，所以 $\C v$ 给出表示 $\varepsilon$。故
\[
\C[G/H]\cong \mathbf1\oplus \varepsilon.
\]

\item \textbf{解答.} 设 $\chi=\chi_{\mathrm{Std}}=(2,0,-1)$。对二维表示，$\mathrm{Sym}^2$ 的特征值是原特征值的二次对称积。

在 $1$ 上，$\mathrm{Std}$ 的特征值为 $1,1$，故
\[
\chi_{\mathrm{Sym}^2(\mathrm{Std})}(1)=1+1+1=3.
\]
在换位上，$\mathrm{Std}$ 的特征值为 $1,-1$，故
\[
\chi_{\mathrm{Sym}^2(\mathrm{Std})}((12))=1+(-1)+1=1.
\]
在三循环上，$\mathrm{Std}$ 的特征值为 $\zeta_3,\zeta_3^2$，故
\[
\chi_{\mathrm{Sym}^2(\mathrm{Std})}((123))
=\zeta_3^2+1+\zeta_3=0.
\]
所以
\[
\chi_{\mathrm{Sym}^2(\mathrm{Std})}=(3,1,0).
\]
分解重数为
\[
\inner{(3,1,0)}{\mathbf1}=\frac16(3+3)=1,\qquad
\inner{(3,1,0)}{\sgn}=\frac16(3-3)=0,
\]
\[
\inner{(3,1,0)}{\mathrm{Std}}=\frac16(3\cdot2)=1.
\]
故
\[
\mathrm{Sym}^2(\mathrm{Std})\cong \mathbf1\oplus \mathrm{Std}.
\]

\item \textbf{解答.} 设 $A\in \End_G(\C[X])$，其矩阵系数为 $a_{xy}$。由 $AG=GA$，若 $(x,y)$ 与 $(x',y')$ 在 $G$ 在 $X\times X$ 上的同一轨道中，则
\[
a_{xy}=a_{x'y'}.
\]
因此 $\End_G(\C[X])$ 的维数等于 $G$ 在 $X\times X$ 上的轨道数。

若作用 $2$-重传递，则 $X\times X$ 恰有两条轨道：
\[
\{(x,x):x\in X\},\qquad \{(x,y):x\neq y\}.
\]
故 $G$-自同态只由这两条轨道上的两个常数决定，从而
\[
\dim \End_G(\C[X])=2.
\]

\item \textbf{解答.} 在基
\[
X^k,\ X^{k-1}Y,\ \dots,\ XY^{k-1},\ Y^k
\]
下，
\[
g=\diag(a,d)
\]
作用为
\[
g\cdot X^{k-i}Y^i=a^{k-i}d^iX^{k-i}Y^i\qquad (0\le i\le k).
\]
故全部特征值为
\[
a^k,\ a^{k-1}d,\ \dots,\ ad^{k-1},\ d^k,
\]
特征标为
\[
\chi_{\mathrm{Sym}^k}(g)=\sum_{i=0}^k a^{k-i}d^i.
\]
若 $a\neq d$，这也可写成
\[
\chi_{\mathrm{Sym}^k}(g)=\frac{a^{k+1}-d^{k+1}}{a-d};
\]
若 $a=d$，则
\[
\chi_{\mathrm{Sym}^k}(g)=(k+1)a^k.
\]

\item \textbf{解答.} 因
\[
|\mathbf P^1(\F_q)|=q+1,
\]
而
\[
\mathrm{St}=\left\{f:\mathbf P^1(\F_q)\to\C:\sum_x f(x)=0\right\}
\]
是一个线性条件定义的子空间，所以
\[
\dim \mathrm{St}=(q+1)-1=q.
\]
再者，任意函数 $f$ 都可唯一写成
\[
f=\left(f-\frac1{q+1}\sum_xf(x)\right)+\frac1{q+1}\sum_xf(x).
\]
第一项坐标和为 $0$，属于 $\mathrm{St}$；第二项为常值函数，属于平凡表示 $\mathbf1$。故
\[
\C[\mathbf P^1(\F_q)]\cong \mathbf1\oplus \mathrm{St}.
\]

\item \textbf{解答.} 定义
\[
\Phi:I(\chi,\chi)\to I(1,1),\qquad \Phi(f)(g)=\chi(\det g)^{-1}f(g).
\]
若
\[
b=\mat{a&*\\0&d}\in B,
\]
则因 $f(bg)=\chi(a)\chi(d)f(g)$，有
\[
\Phi(f)(bg)=\chi(\det(bg))^{-1}f(bg)
=\chi(ad\det g)^{-1}\chi(a)\chi(d)f(g)=\Phi(f)(g).
\]
故 $\Phi(f)\in I(1,1)$。$\Phi$ 显然线性可逆，其逆为
\[
\Psi(\varphi)(g)=\chi(\det g)\varphi(g).
\]
并且对 $x\in G$，
\[
\Phi(x\cdot f)= (\chi\circ\det)(x)\, x\cdot \Phi(f),
\]
故表示意义下
\[
I(\chi,\chi)\cong (\chi\circ\det)\otimes I(1,1).
\]

\item \textbf{解答.} 对 $\dim V$ 作归纳。$\dim V=1$ 时显然。设 $\dim V>1$。若 $A$ 中每个元素都按标量作用，则任取 $0\neq v\in V$ 即为共同特征向量。否则取 $a\in A$ 使其不是标量变换。域代数闭，故 $a$ 有特征值 $\lambda$，对应特征子空间
\[
E=\Ker(a-\lambda I)\neq0.
\]
因 $a$ 非标量，$E\neq V$。对任意 $b\in A$ 与 $v\in E$，
\[
(a-\lambda I)(b\cdot v)=b\cdot(a-\lambda I)(v)=0,
\]
故 $b\cdot v\in E$。于是 $E$ 是非零真 $A$-子模。由归纳假设，$E$ 中存在共同特征向量，也就是 $V$ 中的共同特征向量。

\item \textbf{解答.} 设
\[
f(q)=\sum_{n\ge1}a_nq^n,\qquad T_nf=\lambda_nf.
\]
比较两边 $q^1$ 的系数。由 Hecke 的 $q$-展开公式，
\[
T_nf=\sum_{m\ge1}\left(\sum_{d\mid(m,n)}d^{k-1}a_{mn/d^2}\right)q^m,
\]
故 $q^1$ 的系数为
\[
\sum_{d\mid(1,n)}d^{k-1}a_{n/d^2}=a_n.
\]
而 $\lambda_nf$ 的 $q^1$ 系数为 $\lambda_n a_1=\lambda_n$，因为规范化给出 $a_1=1$。故
\[
\lambda_n=a_n.
\]

\item \textbf{解答.} 由规范化本征模形式满足 $\lambda_n=a_n$，又若 $(m,n)=1$，则
\[
T_mT_n=T_{mn}.
\]
作用于 $f$ 上得
\[
a_ma_n f=\lambda_m\lambda_n f=T_mT_nf=T_{mn}f=\lambda_{mn}f=a_{mn}f.
\]
因 $f\neq0$，故
\[
a_{mn}=a_ma_n.
\]

\item \textbf{解答.} 由 $S_3$ 特征标表
\[
\chi_{\mathbf1}=(1,1,1),\qquad
\chi_{\sgn}=(1,-1,1),\qquad
\chi_{\mathrm{Std}}=(2,0,-1).
\]

首先
\[
\chi_{\sgn\otimes\sgn}=\chi_{\sgn}^2=(1,1,1)=\chi_{\mathbf1},
\]
故
\[
\sgn\otimes\sgn\cong \mathbf1.
\]

其次
\[
\chi_{\sgn\otimes\mathrm{Std}}
=\chi_{\sgn}\chi_{\mathrm{Std}}
=(2,0,-1)=\chi_{\mathrm{Std}},
\]
故
\[
\sgn\otimes\mathrm{Std}\cong \mathrm{Std}.
\]

再者
\[
\chi_{\mathrm{Std}\otimes\mathrm{Std}}
=\chi_{\mathrm{Std}}^2=(4,0,1).
\]
计算内积：
\[
\inner{(4,0,1)}{\mathbf1}=\frac16(4+2)=1,\qquad
\inner{(4,0,1)}{\sgn}=\frac16(4+2)=1,
\]
\[
\inner{(4,0,1)}{\mathrm{Std}}
=\frac16(4\cdot2+2\cdot(-1))=1.
\]
故
\[
\mathrm{Std}\otimes\mathrm{Std}\cong \mathbf1\oplus \sgn\oplus \mathrm{Std}.
\]

\item \textbf{解答.} 将两个复表示分解为不可约直和：
\[
V\cong \bigoplus_i m_iV_i,\qquad W\cong \bigoplus_i n_iV_i,
\]
其中 $V_i$ 两两不等价不可约。则
\[
\chi_V=\sum_i m_i\chi_i,\qquad \chi_W=\sum_i n_i\chi_i.
\]
由第一正交关系，
\[
m_i=\inner{\chi_V}{\chi_i},\qquad n_i=\inner{\chi_W}{\chi_i}.
\]
若 $\chi_V=\chi_W$，则对每个 $i$ 都有 $m_i=n_i$，故两表示具有相同的不可约分解，从而
\[
V\cong W.
\]

\item \textbf{解答.} 对 $D_4$ 的五个共轭类
\[
1,\ r^2,\ r,r^3,\ s,sr^2,\ sr,sr^3
\]
计算顶点不动点数。恒等元固定 $4$ 个顶点；$r^2,r,r^3$ 都不固定顶点；若把正方形顶点取在四个角上，则 $s,sr^2$（过边中点的反射）固定 $0$ 个顶点，而 $sr,sr^3$（过对角线的反射）固定 $2$ 个顶点。故置换特征标为
\[
\chi_{\mathrm{perm}}=(4,0,0,0,2).
\]
与正文中的不可约特征标表作内积：
\[
\inner{\chi_{\mathrm{perm}}}{\chi_{1,1}}=1,\quad
\inner{\chi_{\mathrm{perm}}}{\chi_{1,-1}}=0,\quad
\inner{\chi_{\mathrm{perm}}}{\chi_{-1,1}}=0,
\]
\[
\inner{\chi_{\mathrm{perm}}}{\chi_{-1,-1}}=1,\quad
\inner{\chi_{\mathrm{perm}}}{\rho_{\mathrm{geo}}}=1.
\]
故
\[
\C[\text{四顶点}]\cong \chi_{1,1}\oplus \chi_{-1,-1}\oplus \rho_{\mathrm{geo}}.
\]

\item \textbf{解答.} 将 $V,W$ 分解为不可约直和：
\[
V\cong \bigoplus_i m_iV_i,\qquad W\cong \bigoplus_i n_iV_i.
\]
由 Schur 引理，
\[
\Hom_G(V_i,V_j)=
\begin{cases}
0,& i\neq j,\\
\C,& i=j.
\end{cases}
\]
故
\[
\dim\Hom_G(V,W)=\sum_i m_in_i.
\]
另一方面，
\[
\chi_V=\sum_i m_i\chi_i,\qquad \chi_W=\sum_i n_i\chi_i,
\]
所以
\[
\inner{\chi_V}{\chi_W}
=\sum_{i,j}m_in_j\inner{\chi_i}{\chi_j}
=\sum_i m_in_i.
\]
因此
\[
\dim\Hom_G(V,W)=\inner{\chi_V}{\chi_W}.
\]

\item \textbf{解答.} 定义线性映射
\[
\Phi:V^\vee\otimes W\longrightarrow \Hom_\C(V,W),\qquad
\Phi(\varphi\otimes w)(v)=\varphi(v)w.
\]
它是向量空间同构。检查群作用：
\[
(g\cdot \Phi(\varphi\otimes w))(v)
=g\cdot \Phi(\varphi\otimes w)(g^{-1}v)
=g\cdot \bigl(\varphi(g^{-1}v)w\bigr)
=(g\varphi)(v)\,gw
\]
恰等于
\[
\Phi(g\varphi\otimes gw)(v).
\]
故 $\Phi$ 是 $G$-同构。取不变向量即得
\[
(V^\vee\otimes W)^G\cong \Hom_G(V,W).
\]

与 Schur 引理的关系是：若 $V,W$ 不可约，则
\[
\dim (V^\vee\otimes W)^G=\dim\Hom_G(V,W)=
\begin{cases}
0,& V\not\cong W,\\
1,& V\cong W.
\end{cases}
\]

\item \textbf{解答.} 对任意复表示，总有
\[
\chi_{V^\vee}(g)=\chi_V(g^{-1}).
\]
因为若 $\rho(g)$ 的特征值为 $\lambda_1,\dots,\lambda_n$，则对偶表示上 $\rho^\vee(g)$ 的特征值为 $\lambda_1^{-1},\dots,\lambda_n^{-1}$。又有限群表示可酉化，故每个 $\lambda_i$ 都满足 $|\lambda_i|=1$，于是
\[
\lambda_i^{-1}=\overline{\lambda_i}.
\]
因此
\[
\chi_{V^\vee}(g)=\sum_i\lambda_i^{-1}
=\sum_i\overline{\lambda_i}
=\overline{\sum_i\lambda_i}
=\overline{\chi_V(g)}.
\]

\item \textbf{解答.} 设 $W$ 是 $H$-表示，作用记为 $\sigma$。函数模型下定义
\[
\Ind_H^G W
=\{f:G\to W:\ f(hg)=\sigma(h)f(g)\ \forall\,h\in H,g\in G\},
\]
并令 $G$ 以右平移作用：
\[
(x\cdot f)(g)=f(gx).
\]
取左陪集代表元 $g_1,\dots,g_r$，其中 $r=[G:H]$，使
\[
G=\bigsqcup_{i=1}^r Hg_i.
\]
则任意 $f$ 由 $f(g_1),\dots,f(g_r)$ 唯一决定，因为
\[
f(hg_i)=\sigma(h)f(g_i).
\]
反之任取 $w_i\in W$，可唯一给出这样的 $f$。故作为向量空间，
\[
\Ind_H^G W\cong W^{\oplus r}.
\]
于是
\[
\dim \Ind_H^G W=[G:H]\dim W.
\]

\item \textbf{解答.} Frobenius 互反律的准确表述是：对任意有限群 $G$、子群 $H\subseteq G$、$H$-表示 $W$ 与 $G$-表示 $V$，存在自然同构
\[
\Hom_G(\Ind_H^G W,V)\cong \Hom_H(W,\Res_H^G V).
\]

现设 $H\triangleleft G$ 且 $W=\mathbf1_H$。此时
\[
\Ind_H^G\mathbf1_H
\cong \C[G/H]
\]
是商集上的置换表示。定义
\[
\Theta:\Hom_G(\C[G/H],V)\to V^H,\qquad \Theta(T)=T(e_H).
\]
因 $e_H$ 被 $H$ 固定，对任意 $h\in H$，
\[
h\cdot \Theta(T)=h\cdot T(e_H)=T(h\cdot e_H)=T(e_H)=\Theta(T),
\]
故 $\Theta(T)\in V^H=\Hom_H(\mathbf1_H,\Res_H^G V)$。

反过来，若 $v\in V^H$，定义
\[
T_v(e_{gH})=g\cdot v.
\]
因 $H$ 正规，若 $gH=g'H$，则 $g'=gh$，于是
\[
g'\cdot v=g\cdot(h\cdot v)=g\cdot v,
\]
故定义良好。并且 $T_v$ 显然 $G$-等变。故
\[
\Hom_G(\Ind_H^G\mathbf1_H,V)\cong V^H
=\Hom_H(\mathbf1_H,\Res_H^G V),
\]
这正是该情形下的互反律。

\item \textbf{解答.} 由有限阿贝尔群结构定理，
\[
A\cong C_{n_1}\times \cdots \times C_{n_r}.
\]
又阿贝尔群的不可约复表示全是一维，故 $\widehat A$ 即全部群同态 $A\to\C^\times$ 所成之群。并且
\[
\widehat{A}\cong \widehat{C_{n_1}}\times\cdots\times \widehat{C_{n_r}}.
\]
对每个循环群 $C_n=\langle t\rangle$，取本原 $n$ 次单位根 $\zeta_n$，则
\[
\chi_j(t)=\zeta_n^j\qquad (j\in \Z/n\Z)
\]
给出同构
\[
\widehat{C_n}\cong C_n.
\]
故
\[
\widehat A\cong C_{n_1}\times\cdots\times C_{n_r}\cong A.
\]
此同构依赖于各因子的生成元与单位根选择，故一般非规范。

\item \textbf{解答.} $\GL_2(\F_p)$ 作用在 $\mathbf P^1(\F_p)$ 上。定义
\[
\GL_2(\F_p)/B\longrightarrow \mathbf P^1(\F_p),\qquad gB\mapsto g\cdot[1:0].
\]
因为 $[1:0]$ 的稳定子正是上三角 Borel 子群 $B$，故此映射良定义且为双射。因此
\[
G/B\cong \mathbf P^1(\F_p).
\]
于是
\[
[G:B]=|\mathbf P^1(\F_p)|=p+1.
\]
故主系列表示
\[
I(\chi_1,\chi_2)=\Ind_B^G(\chi_1\boxtimes\chi_2)
\]
的维数为
\[
\dim I(\chi_1,\chi_2)=[G:B]=p+1.
\]

\item \textbf{解答.} 置换表示的特征标总等于不动点个数，因此
\[
\pi(g)=\#\{x\in \mathbf P^1(\F_p):g\cdot x=x\}.
\]
而 $\mathbf P^1(\F_p)$ 的点正是 $\F_p^2$ 中的一维子空间，所以不动点恰是 $g$ 的 $\F_p$-特征直线。

因此：

(1) 若 $g=I$，则所有直线都固定，故
\[
\pi(I)=p+1.
\]

(2) 若 $g$ 为非平凡幺幺元（即 Jordan 形为 $\mat{1&1\\0&1}$，题中“非平凡幂零元”可如此理解），则只有一条特征直线，故
\[
\pi(g)=1.
\]

(3) 若 $g$ 可对角化，则其不动点数等于其在 $\F_p$ 上的特征直线数：
\[
\pi(g)=
\begin{cases}
p+1,& g=\lambda I\ \text{为标量};\\
2,& g\ \text{有两个不同的 }\F_p\text{-特征值};\\
0,& g\ \text{特征多项式在 }\F_p\text{ 上不可约}.
\end{cases}
\]

\item \textbf{解答.} 由
\[
G/B\cong \mathbf P^1(\F_q)
\]
可知
\[
\Ind_B^{\GL_2(\F_q)}\mathbf1 \cong \C[\mathbf P^1(\F_q)].
\]
而正文已证明
\[
\C[\mathbf P^1(\F_q)]\cong \mathbf1\oplus \mathrm{St}.
\]
故
\[
\Ind_B^{\GL_2(\F_q)}\mathbf1\cong \mathbf1\oplus \mathrm{St}.
\]
因此 Steinberg 表示正可定义为诱导表示中除去常值函数那一维平凡部分之后得到的补表示；等价地，
\[
\mathrm{St}\cong \Ind_B^{\GL_2(\F_q)}\mathbf1/\mathbf1.
\]
这就是它最基本的构造方式。

\item \textbf{解答.} 先分轨道写成
\[
X=\bigsqcup_{i=1}^r X_i.
\]
对每条轨道取一点 $x_i$，稳定子为 $H_i$，则
\[
\C[X_i]\cong \Ind_{H_i}^G\mathbf1.
\]
故
\[
\C[X]\cong \bigoplus_{i=1}^r \Ind_{H_i}^G\mathbf1.
\]
平凡表示在 $\C[X]$ 中的重数为
\[
\sum_{i=1}^r \dim\Hom_G(\mathbf1,\Ind_{H_i}^G\mathbf1).
\]
由 Frobenius 互反律，
\[
\Hom_G(\mathbf1,\Ind_{H_i}^G\mathbf1)\cong \Hom_{H_i}(\mathbf1,\mathbf1)\cong \C,
\]
故每条轨道贡献 $1$，总重数即
\[
r=\#(X/G).
\]
若作用传递，则 $r=1$，所以平凡表示恰出现一次。

\item \textbf{解答.} 对规范化 Hecke 本征 cusp form $f$，有
\[
T_nf=a_nf
\]
对一切 $n\ge1$ 成立。将关系
\[
T_pT_{p^r}=T_{p^{r+1}}+p^{k-1}T_{p^{r-1}}
\]
作用于 $f$，得
\[
a_pa_{p^r}f=a_{p^{r+1}}f+p^{k-1}a_{p^{r-1}}f.
\]
因 $f\neq0$，故
\[
a_{p^{r+1}}=a_pa_{p^r}-p^{k-1}a_{p^{r-1}}.
\]

\item \textbf{解答.} 记
\[
A_r=a_{p^r},\qquad A_0=1,\qquad F(X)=\sum_{r\ge0}A_rX^r.
\]
由上一题递推，
\[
A_{r+1}=a_pA_r-p^{k-1}A_{r-1}\qquad (r\ge1).
\]
于是
\[
(1-a_pX+p^{k-1}X^2)F(X)
=1+(A_1-a_pA_0)X+\sum_{r\ge1}(A_{r+1}-a_pA_r+p^{k-1}A_{r-1})X^{r+1}.
\]
因 $A_1=a_p$，且递推式使后面各项全为 $0$，故
\[
(1-a_pX+p^{k-1}X^2)F(X)=1.
\]
因此
\[
\sum_{r\ge0}a_{p^r}X^r=\frac{1}{1-a_pX+p^{k-1}X^2}.
\]

分母是二次式，说明存在两数 $\alpha_p,\beta_p$ 使
\[
1-a_pX+p^{k-1}X^2=(1-\alpha_pX)(1-\beta_pX),
\]
即
\[
\alpha_p+\beta_p=a_p,\qquad \alpha_p\beta_p=p^{k-1}.
\]
这正像某个二维表示在 Frobenius 元上的特征多项式
\[
1-\Tr(\rho(\Frob_p))X+\det(\rho(\Frob_p))X^2.
\]
因此它暗示 Hecke 特征值由某个二维表示控制。

\item \textbf{解答.} 可分三层理解。

\textnormal{(1) 有限维线性代数图景：} 一族可交换算子（这里是 Hecke 算子）作用在有限维空间上时，可寻找共同特征向量。对共同特征向量 $f$，每个 Hecke 算子 $T$ 都满足
\[
Tf=\lambda(T)f,
\]
从而得到 Hecke 代数到标量域的一个代数同态 $T\mapsto\lambda(T)$。

\textnormal{(2) 群表示图景：} Hecke 算子来自双陪集，因此组成某个 Hecke 代数 $\mathcal H(G,K)$。若 $\pi$ 是一个不可约表示，且其 $K$-固定向量空间 $\pi^K$ 为一维，则 $\mathcal H(G,K)$ 在这条线上只能按标量作用；换言之，
\[
\mathcal H(G,K)\to \End(\pi^K)\cong \C
\]
给出一个特征标。Hecke 本征值就是这些标量。

\textnormal{(3) 自守表示图景：} 一个规范化 Hecke 本征 cusp form $f$ 生成一个不可约自守表示
\[
\pi_f=\bigotimes_v' \pi_{f,v}.
\]
对几乎所有素数 $p$，局部分量 $\pi_{f,p}$ 未分歧，故其球向量空间 $\pi_{f,p}^{K_p}$ 为一维。未分歧局部 Hecke 代数
\[
\mathcal H(\GL_2(\Q_p),K_p)
\]
在这条线上按标量作用；该代数特征标给出的标量正是 $f$ 的 Hecke 本征值，特别是标准双陪集对应的标量就是 $a_p(f)$。所以原口号的精确意思是：模形式的 Hecke 本征值，不是孤立数字，而是其所对应自守表示在各未分歧局部 Hecke 代数上的特征标量。
\end{enumerate}
